\documentclass[12pt, a4paper]{report}

% ── Geometry ──
\usepackage[left=1.25in, right=1in, top=1in, bottom=1in]{geometry}

% ── Font: Times New Roman ──
\usepackage{mathptmx}
\usepackage[T1]{fontenc}
\usepackage[utf8]{inputenc}

% ── Spacing ──
\usepackage{setspace}
\onehalfspacing

% ── Page numbers bottom-center ──
\usepackage{fancyhdr}
\pagestyle{fancy}
\fancyhf{}
\fancyfoot[C]{\thepage}
\renewcommand{\headrulewidth}{0pt}
\fancypagestyle{plain}{\fancyhf{}\fancyfoot[C]{\thepage}\renewcommand{\headrulewidth}{0pt}}

% ── Chapter / Section formatting (14pt bold) ──
\usepackage{titlesec}
\titleformat{\chapter}[display]
  {\normalfont\fontsize{14}{17}\bfseries\centering}
  {Chapter-\thechapter}{0pt}{\fontsize{14}{17}\bfseries}
\titlespacing*{\chapter}{0pt}{-20pt}{20pt}

\titleformat{\section}{\normalfont\fontsize{14}{17}\bfseries}{\thesection}{1em}{}
\titleformat{\subsection}{\normalfont\fontsize{14}{17}\bfseries}{\thesubsection}{1em}{}
\titleformat{\subsubsection}{\normalfont\fontsize{14}{17}\bfseries}{\thesubsubsection}{1em}{}

% ── Tables ──
\usepackage{booktabs}
\usepackage{longtable}
\usepackage{array}
\usepackage{tabularx}
\usepackage{multirow}

% ── Misc ──
\usepackage{enumitem}
\usepackage{hyperref}
\hypersetup{colorlinks=true,linkcolor=black,citecolor=black,urlcolor=blue}
\usepackage{amsmath,amssymb}
\usepackage{graphicx}
\usepackage{caption}
\captionsetup{font=normalsize,labelfont=bf}
\usepackage{pifont}   % \ding{51} ✓  \ding{55} ✗
\usepackage{fancyvrb}
\usepackage{textcomp}

\begin{document}
<center>

\begin{center}
{\fontsize{14}{17}\bfseries CAPSTONE PROJECT REPORT}\\
\end{center}
\textbf{(Project Term January--May 2026)}

\&nbsp;

\section{KRISHISAARTHI: AN INTEGRATED AI PLATFORM FOR SMALLHOLDER FARM MANAGEMENT}
\&nbsp;

\subsection{Submitted by}
\begin{table}[htbp]
\centering
\small
\begin{tabularx}{\textwidth}{|X|X|}
\hline
Name & Registration Number \\
\hline
Gautam Chandra Giri & …………………….. \\
Rohan Agarwal & …………………….. \\
Subodh Katana & …………………….. \\
Biswajeet Ray & …………………….. \\
Abhishek Bhardwaj & …………………….. \\
Aditya Pandey & …………………….. \\
\hline
\end{tabularx}
\end{table}

\&nbsp;

\textbf{Project Group Number:} …………

\textbf{Course Code:} ……………………

\&nbsp;

\subsection{Under the Guidance of}
\textbf{(Name of Faculty Mentor with Designation)}

\textbf{School of Computer Science and Engineering}

\textbf{Lovely Professional University, Phagwara, Punjab --- 144411}

</center>

\textbackslash\{\}newpage

<center>

\begin{center}
{\fontsize{14}{17}\bfseries CAPSTONE PROJECT REPORT}\\
\end{center}
\textbf{(Project Term January--May 2026)}

\&nbsp;

\section{KRISHISAARTHI: AN INTEGRATED AI PLATFORM FOR SMALLHOLDER FARM MANAGEMENT}
\&nbsp;

\subsection{Submitted by}
\begin{table}[htbp]
\centering
\small
\begin{tabularx}{\textwidth}{|X|X|}
\hline
Name & Registration Number \\
\hline
Gautam Chandra Giri & …………………….. \\
Rohan Agarwal & …………………….. \\
Subodh Katana & …………………….. \\
Biswajeet Ray & …………………….. \\
Abhishek Bhardwaj & …………………….. \\
Aditya Pandey & …………………….. \\
Anzah Bashir & …………………….. \\
\hline
\end{tabularx}
\end{table}

\&nbsp;

\textbf{Project Group Number:} …………

\textbf{Course Code:} ……………………

\&nbsp;

\subsection{Under the Guidance of}
\textbf{(Name of Faculty Mentor with Designation)}

\textbf{School of Computer Science and Engineering}

\textbf{Lovely Professional University, Phagwara, Punjab --- 144411}

</center>

\textbackslash\{\}newpage

\section{PAC FORM}
*(Attach the PAC form here)*

\textbackslash\{\}newpage

\section{DECLARATION}
We hereby declare that the project work entitled \textbf{"KrishiSaarthi: An Integrated AI Platform for Smallholder Farm Management"} is an authentic record of our own work carried out as requirements of Capstone Project for the award of B.Tech degree in \textbf{Computer Science and Engineering} from Lovely Professional University, Phagwara, under the guidance of \textbf{(Name of Faculty Mentor)}, during January to May 2026. All the information furnished in this capstone project report is based on our own intensive work and is genuine.

\textbf{Project Group Number:} …………

\begin{table}[htbp]
\centering
\small
\begin{tabularx}{\textwidth}{|X|X|}
\hline
Name & Registration Number \\
\hline
Gautam Chandra Giri & …………………….. \\
Rohan Agarwal & …………………….. \\
Subodh Katana & …………………….. \\
Biswajeet Ray & …………………….. \\
Abhishek Bhardwaj & …………………….. \\
Aditya Pandey & …………………….. \\
Anzah Bashir & …………………….. \\
\hline
\end{tabularx}
\end{table}

\&nbsp;

(Signature of Student 1) --- Date:

(Signature of Student 2) --- Date:

(Signature of Student 3) --- Date:

(Signature of Student 4) --- Date:

(Signature of Student 5) --- Date:

(Signature of Student 6) --- Date:

(Signature of Student 7) --- Date:

\textbackslash\{\}newpage

\section{CERTIFICATE}
This is to certify that the declaration statement made by this group of students is correct to the best of my knowledge and belief. They have completed this Capstone Project under my guidance and supervision. The present work is the result of their original investigation, effort and study. No part of the work has ever been submitted for any other degree at any University. The Capstone Project is fit for the submission and partial fulfillment of the conditions for the award of B.Tech degree in \textbf{Computer Science and Engineering} from Lovely Professional University, Phagwara.

\&nbsp;

Signature and Name of the Mentor

Designation

School of Computer Science and Engineering,

Lovely Professional University,

Phagwara, Punjab.

Date:

\textbackslash\{\}newpage

\section{ACKNOWLEDGEMENT}
We would like to express our heartfelt gratitude to everyone who helped us complete this Capstone Project successfully.

First, we sincerely thank our project guide for their valuable guidance, constant motivation, and helpful feedback throughout the development of KrishiSaarthi. Their deep understanding of the subject and technical skills played a big role in shaping both the research direction and the way we built this platform.

We are truly grateful to the Head of the Department of Computer Science and Engineering, and all the faculty members at Lovely Professional University, for giving us the right academic environment and resources to carry out this work.

We also acknowledge the support of the open-source community, especially the developers and maintainers of PyTorch, Django, React, Google Earth Engine, and Google Gemini AI. Their tools and frameworks are the backbone of this platform.

We thank our families and friends for their continuous moral support and patience during the busy phases of research and coding.

Finally, we dedicate this work to the smallholder farming community of India, the people this platform is built for, whose real-world challenges inspired every feature, every design choice, and every line of code in KrishiSaarthi.

\textbackslash\{\}newpage

\section{TABLE OF CONTENTS}
\begin{table}[htbp]
\centering
\small
\begin{tabularx}{\textwidth}{|X|X|X|}
\hline
Section & Title & Page \\
\hline
 & Inner first page & (i) \\
 & PAC form & (ii) \\
 & Declaration & (iii) \\
 & Certificate & (iv) \\
 & Acknowledgement & (v) \\
 & Table of Contents & (vi) \\
\textbf{1} & \textbf{INTRODUCTION} & 1 \\
1.1 & Background and Importance & 1 \\
1.2 & Problem Statement & 4 \\
1.3 & Objectives of the Project & 5 \\
\textbf{2} & \textbf{PROFILE OF THE PROBLEM} & 6 \\
2.1 & Rationale of the Study & 6 \\
2.2 & Scope of the Study & 7 \\
\textbf{3} & \textbf{EXISTING SYSTEM} & 9 \\
3.1 & Introduction & 9 \\
3.2 & Existing Software & 9 \\
3.3 & DFD for Present System & 11 \\
3.4 & What's New in the System to Be Developed & 12 \\
\textbf{4} & \textbf{PROBLEM ANALYSIS} & 14 \\
4.1 & Product Definition & 14 \\
4.2 & Feasibility Analysis & 15 \\
4.3 & Project Plan & 17 \\
\textbf{5} & \textbf{SOFTWARE REQUIREMENT ANALYSIS} & 20 \\
5.1 & Introduction & 20 \\
5.2 & General Description & 20 \\
5.3 & Specific Requirements & 22 \\
\textbf{6} & \textbf{DESIGN} & 26 \\
6.1 & System Design & 26 \\
6.2 & Design Notations & 29 \\
6.3 & Detailed Design & 30 \\
6.4 & Flowcharts & 34 \\
6.5 & Pseudo Code & 37 \\
\textbf{7} & \textbf{TESTING} & 40 \\
7.1 & Functional Testing & 40 \\
7.2 & Structural Testing & 42 \\
7.3 & Levels of Testing & 43 \\
7.4 & Testing the Project & 44 \\
\textbf{8} & \textbf{IMPLEMENTATION} & 47 \\
8.1 & Implementation of the Project & 47 \\
8.2 & Conversion Plan & 49 \\
8.3 & Post-Implementation and Software Maintenance & 50 \\
\textbf{9} & \textbf{PROJECT LEGACY} & 52 \\
9.1 & Current Status of the Project & 52 \\
9.2 & Remaining Areas of Concern & 53 \\
9.3 & Technical and Managerial Lessons Learnt & 54 \\
\textbf{10} & \textbf{USER MANUAL} & 55 \\
\textbf{11} & \textbf{SOURCE CODE / SYSTEM SNAPSHOTS} & 58 \\
\textbf{12} & \textbf{BIBLIOGRAPHY} & 62 \\
\hline
\end{tabularx}
\end{table}

\textbackslash\{\}newpage

\begin{center}
{\fontsize{14}{17}\bfseries 1. INTRODUCTION}\\
\end{center}
\section{1.1. BACKGROUND AND IMPORTANCE}
Agriculture is the backbone of the Indian economy. It employs about 42\% of the country's workforce and contributes around 18.3\% to the GDP (Economic Survey of India, 2024-25). Smallholder farms, which are farms under two hectares, make up roughly 86\% of all farms in India and operate on 47.3\% of total cultivated land (Agricultural Census, 2015-16). Despite their huge role in food security and rural livelihoods, these farmers continue to struggle with a range of connected challenges.

Crop failures from pest attacks and disease outbreaks remain a major threat. Market price swings make planning difficult. Language barriers keep many farmers from getting proper advisory services. And climate change, with its unpredictable weather patterns, is making traditional farming practices less reliable year after year.

Precision agriculture has come up as a promising way to tackle these problems. It uses satellite data, cloud computing, machine learning, and IoT to give field-level guidance to farmers. Commercial platforms like Climate FieldView (by Bayer), CropX, and Plantix have shown what is possible when technology meets farming. But these tools are mostly built for large farms in developed countries. They come with steep licensing costs (USD 500-2,000 per farm per year), need special hardware, or demand advanced technical know-how. That puts them out of reach for the vast majority of smallholder farmers.

However, a few recent advances now make it possible to deliver similar capabilities at almost zero cost:

\begin{enumerate}[leftmargin=*]
  \item \textbf{Free satellite imagery} - Copernicus Sentinel-2 provides cloud-free multispectral images at 10-meter resolution every 5 days, freely available through Google Earth Engine.
  \item \textbf{Lightweight deep learning models} - Architectures like MobileNetV2 (3.4 million parameters, 9.1 MB) can run accurate image classification on regular CPUs, so there is no need for expensive GPUs.
  \item \textbf{Free climate data} - ERA5 reanalysis datasets from ECMWF provide historical and near-real-time climate data at no cost.
  \item \textbf{Large Language Model APIs} - Google Gemini 2.5 Flash offers low-latency, multilingual conversational AI via API, which can power natural language advisory services.
  \item \textbf{Open-source web frameworks} - Django 5.1, React 18, and related libraries give production-grade development abilities without any licensing fees.
\end{enumerate}
This project uses all five of these advances to build a complete, open-source platform that brings precision agriculture to smallholder farmers.

\subsection{1.1.1. The Indian Agricultural Context}
India's farming landscape has some unique features that set it apart:

\textbf{Fragmented land:} The average farm size in India has gone down from 2.28 hectares in 1970-71 to just 1.08 hectares in 2015-16. Over 68\% of all holdings are less than 1 hectare.

\textbf{Diverse climates:} India covers 15 different agro-climatic zones, from the dry deserts of Rajasthan to the humid tropics of Kerala. Any farming AI platform must handle this diversity without needing zone-specific changes.

\textbf{Limited connectivity:} Smartphone use has grown a lot (about 750 million users by 2025), but internet quality is still uneven in rural areas.

\textbf{Language diversity:} India has 22 official languages and over 780 distinct languages. If you only provide advisory services in English, you reach less than 11\% of farmers.

\textbf{Climate vulnerability:} About 52\% of cultivated land in India depends on monsoon rain. More frequent droughts, floods, and unseasonal rainfall directly hit crop yields and farmer incomes.

\section{1.2. PROBLEM STATEMENT}
The central question this project addresses is:

*How can an open-source, AI-powered platform be designed, built, and tested to handle six critical areas of smallholder farm management - crop health monitoring, disease detection, risk forecasting, financial management, carbon credit estimation, and multilingual advisory - all in one integrated system that runs on ordinary hardware without GPU, supports multiple Indian languages, and performs as well as or better than existing commercial tools?*

The specific sub-problems are:

\begin{itemize}[leftmargin=*]
  \item \textbf{SP1:} Can MobileNetV2 (3.4M parameters) match or beat heavyweight models for multi-class crop disease classification, while keeping inference under one second on CPU?
  \item \textbf{SP2:} Can a two-layer LSTM network trained on real-world ERA5 climate data produce reliable daily risk forecasts across different agro-climatic zones?
  \item \textbf{SP3:} Can NDWI time-series analysis automatically detect AWD irrigation cycles and estimate carbon credits using IPCC Tier 2 methodology?
  \item \textbf{SP4:} Can a composite health scoring system that combines CNN, satellite, LSTM, weather, and farming practice signals give a useful crop health assessment?
  \item \textbf{SP5:} Can a multilingual LLM-powered chat interface provide relevant agricultural advice with voice input and sub-3-second response times?
\end{itemize}
\section{1.3. OBJECTIVES OF THE PROJECT}
Three primary objectives guide this work:

\begin{itemize}[leftmargin=*]
  \item \textbf{RO1 - Platform Design:} Design and build KrishiSaarthi as a production-grade platform with a REST API covering all six identified needs, deployed as a containerized application.
  \item \textbf{RO2 - ML Model Development:} Train and test the MobileNetV2 CNN (38-class), LSTM risk model (25 Indian districts), and composite health scoring model against published benchmarks.
  \item \textbf{RO3 - Gap Coverage:} Check platform coverage across all six areas and compare with leading commercial tools and 17 peer-reviewed research papers.
\end{itemize}
\section{1.4. SIGNIFICANCE AND EXPECTED IMPACT}
\begin{enumerate}[leftmargin=*]
  \item \textbf{Making precision agriculture accessible} - Open-source and GPU-free, removing cost and complexity barriers. Total cost drops by over 90\% compared to commercial platforms.
  \item \textbf{New income through carbon credits} - IPCC Tier 2 AWD module can earn between 1,400 and 10,800 INR per AWD cycle while cutting methane emissions by up to 48\%.
  \item \textbf{Breaking the language barrier} - 4 Indian languages plus voice input for the 89.4\% who do not speak English.
  \item \textbf{Advancing research} - Two-stage LLM-augmented inference pipeline is absent from all 17 reviewed papers.
  \item \textbf{Open-source community} - Full code, weights, Docker deployment under MIT license.
\end{enumerate}
\section{1.5. ORGANIZATION OF THE REPORT}
\begin{itemize}[leftmargin=*]
  \item \textbf{Chapter 1} covers introduction, background, problem statement, and objectives.
  \item \textbf{Chapter 2} describes the profile of the problem, rationale, and scope.
  \item \textbf{Chapter 3} reviews the existing systems and what is new in KrishiSaarthi.
  \item \textbf{Chapter 4} presents the problem analysis including feasibility and project plan.
  \item \textbf{Chapter 5} details the software requirement analysis.
  \item \textbf{Chapter 6} covers the system design, flowcharts, and pseudo code.
  \item \textbf{Chapter 7} discusses the testing strategy and results.
  \item \textbf{Chapter 8} describes the implementation, conversion, and maintenance plan.
  \item \textbf{Chapter 9} presents the project legacy.
  \item \textbf{Chapter 10} is the complete user manual.
  \item \textbf{Chapter 11} shows source code highlights and system snapshots.
  \item \textbf{Chapter 12} lists the bibliography.
\end{itemize}
\textbackslash\{\}newpage

\begin{center}
{\fontsize{14}{17}\bfseries 2. PROFILE OF THE PROBLEM}\\
\end{center}
\section{2.1. RATIONALE OF THE STUDY}
The rationale for this project comes from three important observations:

\textbf{1. Unmet needs in a single platform:} Despite significant progress in agricultural AI research, no single tool addresses all the needs of smallholder farmers in one place. Our review of 17 peer-reviewed papers confirmed that no study covers more than three of the eight functional areas we identified. Farmers currently have to juggle multiple apps and services to manage their farms, which is impractical for most.

\textbf{2. Technology is now ready:} The coming together of free satellite data, lightweight ML models, and open-source frameworks now makes it technically and financially possible to deliver enterprise-grade agricultural intelligence at near-zero cost. This was simply not achievable even 3-5 years ago.

\textbf{3. Untapped carbon revenue:} Sustainable practices like AWD irrigation for rice generate measurable carbon credits. The IPCC Tier 2 methodology gives a standard way to quantify this. But no existing platform implements this capability to help farmers earn from sustainable practices.

\subsection{2.1.1. Key Gaps in Current Solutions}
Table 2: Identified Gaps in Existing Smallholder Tooling

\begin{table}[htbp]
\centering
\small
\begin{tabularx}{\textwidth}{|X|X|}
\hline
Gap Area & Description \\
\hline
No integrated platform & Farmers must use 3-5 separate apps for disease detection, weather, finance, and advisory \\
No LLM advisory in Indian languages & Existing chatbots are English-only or rule-based, not AI-powered \\
No financial management & Agricultural AI research ignores the economic side that matters most to farmers \\
No carbon credit tools & No platform translates irrigation practices into monetary carbon credit value \\
Expensive hardware needs & Most solutions need GPUs, IoT sensors, or proprietary devices \\
No explainability & Farmers and extension workers cannot understand why the AI made a prediction \\
\hline
\end{tabularx}
\end{table}

\section{2.2. SCOPE OF THE STUDY}
\subsection{2.2.1. Geographic Scope}
The platform is designed and tested for Indian agricultural conditions. The LSTM risk model is trained on ERA5 climate data covering 25 agricultural districts across different agro-climatic zones:

\begin{itemize}[leftmargin=*]
  \item \textbf{Arid regions:} Jaipur (Rajasthan), Ahmedabad (Gujarat), Jodhpur (Rajasthan)
  \item \textbf{Semi-arid regions:} Hisar (Haryana), Nagpur (Maharashtra), Hyderabad (Telangana)
  \item \textbf{Humid tropical regions:} Ernakulam (Kerala), Coimbatore (Tamil Nadu), Guntur (Andhra Pradesh)
  \item \textbf{Sub-tropical regions:} Ludhiana (Punjab), Lucknow (Uttar Pradesh), Patna (Bihar)
  \item \textbf{Highland regions:} Shimla (Himachal Pradesh), Imphal (Manipur)
\end{itemize}
\subsection{2.2.2. Crop Coverage}
The MobileNetV2 disease classifier handles 38 classes across 14 crop species from the PlantVillage dataset: tomato (10 classes), grape (4), apple (4), corn (4), potato (3), pepper (2), strawberry (2), cherry (2), peach (2), citrus (1), blueberry (1), raspberry (1), soybean (1), and squash (1).

\subsection{2.2.3. Functional Scope}
The platform covers six functional areas end-to-end:

\begin{enumerate}[leftmargin=*]
  \item Satellite-based crop monitoring (NDVI, EVI, SAVI, NDWI via Google Earth Engine)
  \item AI-powered disease detection (38-class CNN with LLM pre-validation)
  \item Temporal risk forecasting (LSTM-based daily risk prediction for 25 districts)
  \item Financial management with P\&L analysis (costs, revenue, seasons, market prices, insurance)
  \item Carbon credit estimation (IPCC Tier 2 AWD methodology)
  \item Multilingual conversational advisory (4 languages plus voice input)
\end{enumerate}
\subsection{2.2.4. Exclusions}
The project does not include real-world field trials with actual farmers (planned as future work), IoT sensor integration, drone-based imagery analysis, or blockchain-based supply chain tracking. The disease model is tested on PlantVillage laboratory data; field-level validation is planned for the future.

\textbackslash\{\}newpage

\begin{center}
{\fontsize{14}{17}\bfseries 3. EXISTING SYSTEM}\\
\end{center}
\section{3.1. INTRODUCTION}
Before building KrishiSaarthi, we studied the existing tools and platforms available to Indian farmers. This chapter looks at what is already out there, how those systems work, what they lack, and what new things our system brings to the table.

The agricultural technology space has grown rapidly in recent years. Several platforms now offer disease detection, satellite monitoring, or weather information. However, most of these are either too expensive for small farmers, limited in what they do, or designed for large commercial farms in developed countries. None of them brings together all six functional areas (crop health, disease detection, risk forecasting, finance, carbon credits, and multilingual advisory) under one roof.

\section{3.2. EXISTING SOFTWARE}
\subsection{3.2.1. Climate FieldView (Bayer)}
Climate FieldView is one of the most well-known precision agriculture platforms. It offers satellite-based field health mapping, yield prediction, and seed selection advice. However, it costs USD 500-2,000 per farm per year, needs specialized hardware, works only in English, and is built for large US farms. Indian smallholders with 1-2 hectares cannot justify this expense.

\subsection{3.2.2. CropX}
CropX focuses on soil moisture sensing and irrigation management. It provides good irrigation advice, but it requires proprietary IoT sensors that cost USD 200 or more per unit. It covers a limited number of crops, has no disease detection, and does not help with financial management.

\subsection{3.2.3. Plantix}
Plantix is a mobile app for crop disease identification using photos. It has a large community of users and covers many crops. However, it is limited to just disease identification. It has no satellite integration, no financial tools, no carbon credit features, and uses a freemium model that restricts access to advanced features.

\subsection{3.2.4. Kisan Suvidha (Government of India)}
Kisan Suvidha is a government app that provides market prices, weather updates, and dealer info. It is free but has no AI or ML capabilities. The interface is basic, the data is often static, and it simply aggregates information without any personalized analysis.

\subsection{3.2.5. eNAM (Electronic National Agriculture Market)}
eNAM is a government electronic marketplace for agricultural produce. It helps with price discovery and market transactions but has no farm management features, no crop health monitoring, and no advisory services.

\subsection{3.2.6. Comparison of Existing Platforms}
Table 3: Comparison of Existing Platforms with KrishiSaarthi

\begin{table}[htbp]
\centering
\small
\begin{tabularx}{\textwidth}{|X|X|X|X|X|X|}
\hline
Feature & Climate FieldView & CropX & Plantix & Kisan Suvidha & KrishiSaarthi \\
\hline
Disease Detection & No & No & Yes (limited) & No & Yes (38-class CNN + LLM) \\
Satellite Monitoring & Yes & No & No & No & Yes (4 indices) \\
Risk Prediction & Yes & No & No & No & Yes (LSTM, 25 districts) \\
Financial Management & No & No & No & No & Yes (Full P\&L) \\
Carbon Credits & No & No & No & No & Yes (IPCC Tier 2) \\
Multilingual Support & No & No & Partial & Hindi/English & Yes (4 languages + voice) \\
AI Chat Advisory & No & No & No & No & Yes (Gemini 2.5 Flash) \\
Open Source & No & No & No & No & Yes (MIT License) \\
Annual Cost & \$500-\$2,000 & \$200+ sensors & Freemium & Free & Free \\
\hline
\end{tabularx}
\end{table}

\section{3.3. DFD FOR PRESENT SYSTEM}
The current farm management process without a unified platform works like this:

\begin{Verbatim}[fontsize=\small]
Level 0 DFD - Current Situation (Without KrishiSaarthi)

                +----------------+
                |    Farmer      |
                +-------+--------+
                        |
        +---------------+---------------+-----------------+
        v               v               v                 v
  +-----------+   +-----------+   +--------------+   +-----------+
  | Plantix   |   | Weather   |   | Manual       |   | Mandi     |
  | (Disease  |   | App       |   | Record       |   | Visit     |
  | only)     |   | (Weather  |   | Keeping      |   | (Prices)  |
  |           |   | only)     |   | (Pen/Paper)  |   |           |
  +-----------+   +-----------+   +--------------+   +-----------+

  Problems:
  - No unified view of farm status
  - No risk prediction or early warning
  - No satellite data for crop monitoring
  - No AI advisory in local language
  - Data is scattered across apps and paper
  - No carbon credit tracking
  - Language barriers in most apps
\end{Verbatim}
In the current system, a farmer has to use different apps for different things. Disease checking happens through one app, weather through another, financial records on paper, and market prices require a physical visit to the mandi. There is no connection between these tools, and the farmer gets no integrated insight.

\section{3.4. WHAT'S NEW IN THE SYSTEM TO BE DEVELOPED}
KrishiSaarthi brings several entirely new capabilities that no existing platform offers:

\textbf{1. All-in-one integrated platform:} Instead of juggling 4-5 separate apps and manual registers, a farmer gets everything in one place. Disease detection, satellite monitoring, risk alerts, financial management, crop planning, and AI advice are all available from a single dashboard.

\textbf{2. Two-stage LLM-augmented disease detection:} Before running the CNN, images are first validated by Gemini Vision to check if the uploaded photo is actually a plant leaf. This prevents false predictions on non-plant images (like a selfie or a landscape photo), which is a common real-world problem that no other platform handles.

\textbf{3. IPCC Tier 2 carbon credit estimation:} This is the first platform to calculate carbon credits from AWD irrigation practices using internationally accepted IPCC methodology. Farmers can see how much money they could earn by adopting sustainable water management.

\textbf{4. Composite health scoring:} Instead of a single metric, KrishiSaarthi fuses five different signals (CNN disease probability, satellite NDVI, LSTM risk, weather conditions, and farming practices) into one health score from 0 to 100. This gives a much more complete picture of crop health than any individual model can provide.

\textbf{5. Multilingual AI advisory with voice:} The Gemini-powered chatbot understands and responds in English, Hindi, Punjabi, and Malayalam. It also accepts voice input through the Web Speech API, which helps farmers who are not comfortable typing.

\textbf{6. Explainable AI:} Both the CNN and LSTM models come with explainability tools. Grad-CAM++ shows which parts of the leaf image the CNN focused on, and SHAP analysis shows which climate factors matter most for risk predictions. This builds trust in the AI system.

\textbackslash\{\}newpage

\begin{center}
{\fontsize{14}{17}\bfseries 4. PROBLEM ANALYSIS}\\
\end{center}
\section{4.1. PRODUCT DEFINITION}
KrishiSaarthi (Sanskrit for "Agriculture Companion") is an open-source, AI-powered web platform designed specifically for smallholder farmers in India. It brings together six critical farm management capabilities into one unified system.

\subsection{4.1.1. Product Overview}
The platform is built as a full-stack web application with these key components:

\begin{itemize}[leftmargin=*]
  \item \textbf{Frontend:} React 18 with TypeScript, featuring 79 components organized across 21 dashboard modules. It uses TailwindCSS 4 with Shadcn UI for a clean, modern interface. Leaflet.js provides interactive maps for field boundary drawing, and Recharts handles all data visualization.
  \item \textbf{Backend:} Django 5.1 with Django REST Framework, exposing 45 REST endpoints across 8 functional modules. It uses PostgreSQL in production and SQLite for development. The API is fully documented with OpenAPI 3.0 via drf-spectacular.
  \item \textbf{ML/AI Layer:} A MobileNetV2 CNN for 38-class disease classification (9.1 MB model), a two-layer LSTM for risk prediction (208 KB model), and Google Gemini 2.5 Flash for image validation and conversational advisory.
  \item \textbf{External Services:} Google Earth Engine for satellite vegetation indices, OpenWeather API for weather data, and ERA5 reanalysis data for LSTM training.
\end{itemize}
\subsection{4.1.2. Target Users}
The primary users are smallholder farmers in India who:

\begin{itemize}[leftmargin=*]
  \item Own or operate farms under 2 hectares
  \item Have basic smartphone and internet access
  \item May not be fluent in English
  \item Cannot afford commercial precision agriculture tools
  \item Need help with both agronomic and financial decision-making
\end{itemize}
Secondary users include agricultural extension workers, NGOs working with farming communities, and agricultural researchers.

\section{4.2. FEASIBILITY ANALYSIS}
\subsection{4.2.1. Technical Feasibility}
The project is technically feasible because:

\begin{itemize}[leftmargin=*]
  \item \textbf{ML models run on CPU:} MobileNetV2 (3.4M parameters) achieves sub-millisecond inference on CPU. No GPU is needed. The LSTM model (51,265 parameters) is even lighter.
  \item \textbf{Free data sources:} Sentinel-2 imagery via Earth Engine, ERA5 climate data via Open-Meteo, and OpenWeather API all provide free tiers sufficient for the platform's needs.
  \item \textbf{Mature frameworks:} Django 5.1 and React 18 are battle-tested frameworks with large ecosystems, excellent documentation, and active communities.
  \item \textbf{Cloud deployment ready:} Docker Compose configurations for both development and production are provided. The entire stack runs on a single server with 2 GB RAM.
\end{itemize}
Table 4: Technology Stack Summary

\begin{table}[htbp]
\centering
\small
\begin{tabularx}{\textwidth}{|X|X|X|X|}
\hline
Layer & Technology & Version & Purpose \\
\hline
Frontend & React + TypeScript & 18.x / 5.0+ & Component-based SPA \\
Styling & TailwindCSS + Shadcn UI & 4.0 & Modern UI components \\
Maps & Leaflet.js & 1.9+ & Interactive field maps \\
Charts & Recharts & 2.x & Data visualization \\
Backend & Django + DRF & 5.1 & REST API (45 endpoints) \\
Deep Learning & PyTorch & 2.0+ & CNN and LSTM inference \\
LLM & Google Gemini & 2.5 Flash & Vision + Chat advisory \\
Satellite & Google Earth Engine & Python API & Vegetation indices \\
Database & PostgreSQL / SQLite & 14+ / 3.x & Data storage \\
Cache & Redis & 7.x & Session and response cache \\
Containers & Docker + Compose & 24.x & Deployment \\
Web Server & Nginx + Gunicorn & latest / 21.x & Production serving \\
Monitoring & Prometheus + Grafana & latest & Observability \\
CI/CD & GitHub Actions & - & Automated pipeline \\
\hline
\end{tabularx}
\end{table}

\subsection{4.2.2. Economic Feasibility}
The project is economically viable because:

\begin{itemize}[leftmargin=*]
  \item \textbf{Zero licensing costs:} All frameworks, libraries, and data sources used are either open-source or provide free tiers.
  \item \textbf{Commodity hardware:} The platform runs on a standard server or VPS costing as little as USD 5-20/month. No GPU instances are needed.
  \item \textbf{Free external APIs:} Google Earth Engine, ERA5, and Gemini API free tiers handle the expected load for a smallholder-focused deployment.
  \item \textbf{Comparison with alternatives:} Commercial platforms cost USD 500-2,000 per farm per year. KrishiSaarthi reduces this to near-zero.
\end{itemize}
\subsection{4.2.3. Operational Feasibility}
The platform is operationally feasible because:

\begin{itemize}[leftmargin=*]
  \item \textbf{Web-based access:} No app installation needed. Works on any device with a browser.
  \item \textbf{Multilingual interface:} Four Indian languages reduce the language barrier.
  \item \textbf{Voice input:} Farmers who cannot type comfortably can speak their queries.
  \item \textbf{Responsive design:} The UI works on both desktop and mobile screens.
  \item \textbf{Offline consideration:} The system is designed to degrade gracefully when connectivity is poor. Cached data and fallback values are used when external APIs are unavailable.
\end{itemize}
\section{4.3. PROJECT PLAN}
\subsection{4.3.1. Phase-wise Development Plan}
The project was executed across four phases spanning the academic year 2025-26:

Table 5: Phase-wise Development Timeline

\begin{table}[htbp]
\centering
\small
\begin{tabularx}{\textwidth}{|X|X|X|X|}
\hline
Phase & Duration & Activities & Deliverables \\
\hline
Phase 1: Research and Planning & Aug-Sep 2025 (8 weeks) & Literature review, requirements analysis, architecture design, technology selection & Literature survey, problem statement, architecture document, approved proposal \\
Phase 2: Core Development & Oct-Dec 2025 (12 weeks) & Backend API, database design, ML model training (CNN, LSTM), satellite integration, frontend scaffold & Working backend with 45 endpoints, trained ML models, Earth Engine integration \\
Phase 3: Integration and Features & Jan-Feb 2026 (8 weeks) & Dashboard development, chat integration, financial modules, planning modules, carbon credit engine, multilingual support & Complete full-stack application, all 21 dashboard modules \\
Phase 4: Testing and Documentation & Mar-Apr 2026 (8 weeks) & Performance evaluation, security audit, explainability analysis, Docker deployment, report writing & Evaluation results, Docker images, capstone report \\
\hline
\end{tabularx}
\end{table}

\subsection{4.3.2. Gantt Chart}
\begin{Verbatim}[fontsize=\small]
Phase / Task                    Aug  Sep  Oct  Nov  Dec  Jan  Feb  Mar  Apr
                                 |    |    |    |    |    |    |    |    |
Phase 1: Research & Planning    ========
  Literature Review             ====
  Requirements Analysis              ====
  Architecture Design                ====
  Project Proposal                   ====

Phase 2: Core Development                ====================
  Django Backend Setup                   ====
  Database Schema Design                 ====
  REST API Development                   ============
  CNN Model Training                          ========
  LSTM Model Training                         ========
  Earth Engine Integration                         ========
  Frontend Scaffold                                ========

Phase 3: Integration & Features                          ============
  Dashboard Modules (21)                                  ========
  Financial Management                                    ========
  Planning Modules                                        ========
  Chat/LLM Integration                                  ====
  Carbon Credit Engine                                        ====
  Multilingual + Voice                                        ====

Phase 4: Testing & Docs                                          ============
  Performance Evaluation                                          ====
  Security Audit                                                  ====
  Docker Deployment                                                    ====
  Report Writing                                                  ========
  Final Presentation                                                   ====
\end{Verbatim}
\subsection{4.3.3. Milestone Tracking}
Table 6: Milestone Tracking Table

\begin{table}[htbp]
\centering
\small
\begin{tabularx}{\textwidth}{|X|X|X|X|}
\hline
Milestone & Target Date & Actual Date & Status \\
\hline
M1: Literature review complete & Sep 30, 2025 & Sep 28, 2025 & Completed \\
M2: Project proposal approved & Oct 15, 2025 & Oct 12, 2025 & Completed \\
M3: Backend API (v1) functional & Nov 30, 2025 & Dec 5, 2025 & Completed (5 days late) \\
M4: CNN model trained (>=95\% accuracy) & Dec 15, 2025 & Dec 10, 2025 & Completed (99.03\%) \\
M5: LSTM model trained (>=90\% accuracy) & Dec 31, 2025 & Dec 22, 2025 & Completed (97.56\%) \\
M6: Earth Engine integration working & Jan 15, 2026 & Jan 10, 2026 & Completed \\
M7: All 21 dashboard modules functional & Feb 15, 2026 & Feb 18, 2026 & Completed (3 days late) \\
M8: Multilingual support (4 languages) & Feb 28, 2026 & Feb 25, 2026 & Completed \\
M9: Docker deployment validated & Mar 15, 2026 & Mar 12, 2026 & Completed \\
M10: Performance evaluation complete & Mar 31, 2026 & Apr 3, 2026 & Completed \\
M11: Final report submitted & Apr 15, 2026 & Apr 11, 2026 & Completed \\
\hline
\end{tabularx}
\end{table}

\subsection{4.3.4. Risk Mitigation Plan}
Table 7: Risk Mitigation Matrix

\begin{table}[htbp]
\centering
\small
\begin{tabularx}{\textwidth}{|X|X|X|X|X|}
\hline
Risk & Probability & Impact & Mitigation Strategy & Outcome \\
\hline
Earth Engine API quota exhaustion & Medium & High & Circuit breaker pattern with deterministic NDVI fallback & Implemented successfully \\
CNN model accuracy below target & Low & Critical & Progressive unfreezing strategy; data augmentation & Exceeded target (99.03\% vs 95\%) \\
Gemini API rate limits on chat & Medium & Medium & Response caching; fallback to rule-based advisory & Rate limits not hit during testing \\
Database performance at scale & Low & High & Indexing strategy; query optimization; pagination & Sub-15ms single-field queries achieved \\
Team member availability conflicts & Medium & Medium & Modular architecture enabling parallel development & No critical delays \\
Security vulnerabilities & Medium & Critical & Comprehensive security audit; automated scanning & 42 issues identified and prioritized \\
\hline
\end{tabularx}
\end{table}

\textbackslash\{\}newpage

\begin{center}
{\fontsize{14}{17}\bfseries 5. SOFTWARE REQUIREMENT ANALYSIS}\\
\end{center}
\section{5.1. INTRODUCTION}
This chapter describes the software requirements for the KrishiSaarthi platform. The requirements were gathered through a combination of literature review, analysis of existing platforms, understanding of smallholder farmer needs, and technical feasibility assessment. The requirements are organized following the IEEE 830 standard for Software Requirements Specifications.

The primary goal is to build a system that is technically strong but also simple enough for farmers with limited digital literacy to use. Every feature was evaluated against the question: "Will a farmer with a basic smartphone and limited English skills be able to use this?"

\section{5.2. GENERAL DESCRIPTION}
\subsection{5.2.1. Product Perspective}
KrishiSaarthi is a standalone web application that integrates with several external services. It is not a replacement for any single existing tool but rather a unified platform that combines the functionality of multiple separate applications into one cohesive system.

The system interacts with:

\begin{itemize}[leftmargin=*]
  \item \textbf{Google Earth Engine} for satellite imagery and vegetation index computation
  \item \textbf{OpenWeather API} for real-time weather data and forecasts
  \item \textbf{Google Gemini 2.5 Flash} for image validation and conversational AI
  \item \textbf{ERA5 Reanalysis} (via Open-Meteo) for historical climate data used in model training
\end{itemize}
\subsection{5.2.2. Product Functions}
The platform provides the following major functions:

\begin{enumerate}[leftmargin=*]
  \item \textbf{User Authentication and Profile Management} - Secure registration, login, password reset, and profile management with token-based authentication.
  \item \textbf{Field Management} - Create, view, update, and delete agricultural fields. Each field is defined by a name, crop type, and polygon boundary drawn on an interactive map.
  \item \textbf{Satellite Vegetation Analysis} - On-demand computation of NDVI, EVI, SAVI, and NDWI indices for any registered field using Sentinel-2 satellite imagery via Google Earth Engine.
  \item \textbf{Crop Disease Detection} - Upload a crop leaf image to get a 38-class disease classification with confidence scores and top-3 predictions. Includes Gemini Vision pre-validation to reject non-plant images.
  \item \textbf{Risk Prediction} - LSTM-based daily risk assessment using time-series climate data. Provides risk level (High/Medium/Low) and actionable recommendations.
  \item \textbf{Composite Health Score} - A single 0-100 score that fuses CNN, satellite, LSTM, weather, and practice signals to give an overall crop health assessment.
  \item \textbf{Financial Management} - Track costs (seeds, fertilizer, labor, equipment), record revenue (sales, crop prices), manage seasons, view P\&L dashboards, check market prices, discover government schemes, and manage crop insurance claims.
  \item \textbf{Resource Planning} - Season calendar management, inventory tracking, labor management, equipment scheduling, and crop rotation planning.
  \item \textbf{AWD Detection and Carbon Credits} - Automatic detection of AWD irrigation patterns from NDWI time-series, with IPCC Tier 2 carbon credit estimation.
  \item \textbf{AI Chat Advisory} - Gemini-powered conversational interface with multilingual support (English, Hindi, Punjabi, Malayalam) and voice input.
\end{enumerate}
\subsection{5.2.3. User Characteristics}
The primary users are Indian smallholder farmers who:

\begin{itemize}[leftmargin=*]
  \item May have limited formal education (Class 5 to Class 12)
  \item Primarily speak Hindi, Punjabi, Malayalam, or other regional languages
  \item Have access to smartphones with intermittent internet connectivity
  \item Are not familiar with complex software interfaces
  \item Need practical, actionable farming advice, not academic information
\end{itemize}
\subsection{5.2.4. Constraints}
\begin{itemize}[leftmargin=*]
  \item The system must work without GPU hardware
  \item External API calls (Earth Engine, Gemini) may be slow or unavailable
  \item The system must handle intermittent internet connectivity gracefully
  \item The disease detection model is currently limited to 14 crop species from PlantVillage
  \item Voice input depends on browser support for the Web Speech API
\end{itemize}
\section{5.3. SPECIFIC REQUIREMENTS}
\subsection{5.3.1. Functional Requirements}
\textbf{FR-01: User Registration and Authentication}

\begin{itemize}[leftmargin=*]
  \item The system shall allow users to register with username, email, and password.
  \item The system shall authenticate users using token-based authentication (RFC 6750).
  \item The system shall support password reset via email.
  \item The system shall enforce rate limiting (5 login attempts per minute per IP).
\end{itemize}
\textbf{FR-02: Field Management}

\begin{itemize}[leftmargin=*]
  \item The system shall allow users to create fields by drawing polygon boundaries on a map.
  \item The system shall store field data including name, crop type, area, and GeoJSON polygon.
  \item The system shall allow users to view, edit, and delete their fields.
  \item The system shall ensure users can only access their own fields (user-scoped queries).
\end{itemize}
\textbf{FR-03: Satellite Vegetation Analysis}

\begin{itemize}[leftmargin=*]
  \item The system shall compute NDVI, EVI, SAVI, and NDWI for any registered field.
  \item The system shall use Sentinel-2 L2A imagery via Google Earth Engine.
  \item The system shall apply cloud masking using QA60 band bits.
  \item The system shall provide fallback NDVI estimation when Earth Engine is unavailable.
\end{itemize}
\textbf{FR-04: Disease Detection}

\begin{itemize}[leftmargin=*]
  \item The system shall accept crop leaf images (JPEG/PNG, up to 10 MB).
  \item The system shall validate images using Gemini Vision before CNN inference.
  \item The system shall classify images into 38 PlantVillage disease classes.
  \item The system shall return predicted class, confidence score, and top-3 predictions.
\end{itemize}
\textbf{FR-05: Risk Prediction}

\begin{itemize}[leftmargin=*]
  \item The system shall predict daily crop risk using the trained LSTM model.
  \item The system shall accept time-series input of NDVI, rainfall, temperature, and soil moisture.
  \item The system shall return risk probability, risk level, and recommendations.
\end{itemize}
\textbf{FR-06: Health Score}

\begin{itemize}[leftmargin=*]
  \item The system shall compute a composite health score (0-100) for each field.
  \item The system shall fuse CNN, NDVI, LSTM, weather, and practice signals.
  \item The system shall assign a rating: Excellent, Good, Fair, Poor, or Critical.
\end{itemize}
\textbf{FR-07: Financial Management}

\begin{itemize}[leftmargin=*]
  \item The system shall allow users to record costs by category (seeds, fertilizer, labor, equipment, other).
  \item The system shall allow users to record revenue entries.
  \item The system shall compute profit/loss summaries by season.
  \item The system shall display current market prices for common commodities.
  \item The system shall list relevant government schemes and subsidies.
  \item The system shall support crop insurance claim management.
\end{itemize}
\textbf{FR-08: Resource Planning}

\begin{itemize}[leftmargin=*]
  \item The system shall support season calendar management.
  \item The system shall track inventory items (seeds, fertilizers, pesticides, tools).
  \item The system shall manage labor records (workers, tasks, hours, wages).
  \item The system shall schedule equipment bookings.
  \item The system shall recommend crop rotation plans.
\end{itemize}
\textbf{FR-09: AWD and Carbon Credits}

\begin{itemize}[leftmargin=*]
  \item The system shall detect AWD irrigation patterns from NDWI time-series.
  \item The system shall calculate carbon credits using IPCC Tier 2 methodology.
  \item The system shall estimate monetary value of carbon credits in both USD and INR.
\end{itemize}
\textbf{FR-10: AI Chat Advisory}

\begin{itemize}[leftmargin=*]
  \item The system shall provide AI-powered agricultural advice using Google Gemini.
  \item The system shall support four languages: English, Hindi, Punjabi, Malayalam.
  \item The system shall inject field-specific context (NDVI, weather, risk) into queries.
  \item The system shall maintain conversation history for multi-turn dialogue.
  \item The system shall support voice input via Web Speech API.
\end{itemize}
\subsection{5.3.2. Non-Functional Requirements}
\textbf{NFR-01: Performance}

\begin{itemize}[leftmargin=*]
  \item Database-only endpoints shall respond in under 50 ms.
  \item ML inference endpoints shall respond in under 2 seconds.
  \item Chat advisory shall respond in under 3 seconds.
  \item The system shall support at least 10 concurrent users on a single server.
\end{itemize}
\textbf{NFR-02: Security}

\begin{itemize}[leftmargin=*]
  \item All passwords shall be hashed using PBKDF2-SHA256 (600,000 iterations).
  \item All API endpoints (except health probes) shall require authentication.
  \item PII data (bank accounts, IFSC codes) shall be encrypted at rest using Fernet.
  \item The system shall prevent CSRF, XSS, and SQL injection attacks.
\end{itemize}
\textbf{NFR-03: Reliability}

\begin{itemize}[leftmargin=*]
  \item Circuit breaker patterns shall protect against external API failures.
  \item Fallback values shall be provided when models or APIs are unavailable.
  \item Health probes (/health/live and /health/ready) shall support container orchestration.
\end{itemize}
\textbf{NFR-04: Usability}

\begin{itemize}[leftmargin=*]
  \item The interface shall be responsive and work on screens from 320px to 1920px width.
  \item The system shall support four Indian languages without page reload.
  \item Error messages shall be clear, actionable, and non-technical.
\end{itemize}
\textbf{NFR-05: Maintainability}

\begin{itemize}[leftmargin=*]
  \item The codebase shall follow modular architecture with clear separation of concerns.
  \item All API endpoints shall be documented with OpenAPI 3.0 (Swagger/ReDoc).
  \item Docker Compose configurations shall be provided for both development and production.
\end{itemize}
\subsection{5.3.3. Hardware Requirements}
Table 8: Hardware Requirements

\begin{table}[htbp]
\centering
\small
\begin{tabularx}{\textwidth}{|X|X|X|}
\hline
Component & Minimum & Recommended \\
\hline
CPU & 2 cores (x86\_64) & 4 cores \\
RAM & 2 GB & 4 GB \\
Storage & 5 GB & 10 GB \\
GPU & Not required & Not required \\
Network & Broadband internet & Broadband internet \\
\hline
\end{tabularx}
\end{table}

\subsection{5.3.4. Software Requirements}
Table 9: Software Requirements

\begin{table}[htbp]
\centering
\small
\begin{tabularx}{\textwidth}{|X|X|}
\hline
Component & Requirement \\
\hline
Operating System & Linux (Ubuntu 22.04+), Windows 10+, or macOS 12+ \\
Python & 3.11 or higher \\
Node.js & 18 or higher \\
Database & PostgreSQL 14+ (production) or SQLite 3 (development) \\
Web Browser & Chrome 90+, Firefox 90+, Safari 15+, Edge 90+ \\
Docker & 24.x (optional, for containerized deployment) \\
\hline
\end{tabularx}
\end{table}

\textbackslash\{\}newpage

\begin{center}
{\fontsize{14}{17}\bfseries 6. DESIGN}\\
\end{center}
\section{6.1. SYSTEM DESIGN}
\subsection{6.1.1. Architecture Overview}
KrishiSaarthi follows a four-tier layered architecture that separates concerns and allows each layer to scale independently:

\begin{Verbatim}[fontsize=\small]
+---------------------------------------------------------------------+
|                       CLIENT LAYER (Tier 1)                          |
|  +---------------------------------------------------------------+  |
|  |  React 18 + TypeScript + TailwindCSS 4 + Vite                  |  |
|  |  - 21 Dashboard Features    - Real-time Charts (Recharts)      |  |
|  |  - Interactive Maps (Leaflet) - Multi-language (i18n Hook)     |  |
|  |  - Voice Input (Web Speech)  - 79 UI Components               |  |
|  +---------------------------------------------------------------+  |
+---------------------------------------------------------------------+
                               | REST API (45 endpoints)
                               v
+---------------------------------------------------------------------+
|                      API GATEWAY (Tier 2)                            |
|  +---------------------------------------------------------------+  |
|  |  Django 5.1 + Django REST Framework                            |  |
|  |  - Token Authentication      - CORS Configuration             |  |
|  |  - Rate Limiting (5/min)     - Request Validation (DRF)       |  |
|  |  - CSRF Protection           - Structured Error Responses     |  |
|  +---------------------------------------------------------------+  |
+---------------------------------------------------------------------+
                               |
         +---------------------+---------------------+
         v                     v                     v
+-----------------+   +-----------------+   +-----------------+
|  FIELD MODULE   |   | FINANCE MODULE  |   | PLANNING MODULE |
| - EE Analysis   |   | - Cost Tracking |   | - Calendar      |
| - Health Score  |   | - Revenue       |   | - Inventory     |
| - Disease Detect|   | - P&L Dashboard |   | - Labor         |
| - Irrigation    |   | - Market Prices |   | - Equipment     |
| - AWD / Carbon  |   | - Gov. Schemes  |   | - Crop Rotation |
| - Alerts        |   | - Insurance     |   |                 |
+-----------------+   +-----------------+   +-----------------+
         |                     |                     |
         +---------------------+---------------------+
                               v
+---------------------------------------------------------------------+
|                   ML/AI INFERENCE LAYER (Tier 4)                     |
|  +--------------+  +--------------+  +--------------------------+   |
|  | MobileNetV2  |  | LSTM Risk    |  | Google Gemini 2.5 Flash  |   |
|  | Disease CNN  |  | Prediction   |  | - Image Pre-validation   |   |
|  | 9.1 MB Model |  | 208 KB Model |  | - Chat Advisory          |   |
|  | 38 Classes   |  | 51,265 Params|  | - 4 Languages            |   |
|  +--------------+  +--------------+  +--------------------------+   |
+---------------------------------------------------------------------+
                               |
                               v
+---------------------------------------------------------------------+
|                     EXTERNAL SERVICES                                |
|  +--------------+  +--------------+  +--------------------------+   |
|  | Google Earth |  | OpenWeather  |  | ERA5 Climate Reanalysis  |   |
|  | Engine (GEE) |  | API          |  | (Training Data)          |   |
|  +--------------+  +--------------+  +--------------------------+   |
+---------------------------------------------------------------------+
\end{Verbatim}
Figure 1: Four-Tier System Architecture Diagram

\subsection{6.1.2. Domain Service Modules}
Table 10: Domain Service Modules and Endpoints

\begin{table}[htbp]
\centering
\small
\begin{tabularx}{\textwidth}{|X|X|X|}
\hline
Module & Endpoints & Key Capabilities \\
\hline
Authentication & 5 & Token auth, registration, password reset \\
Field & 8 & Earth Engine analysis, health score, disease detection, AWD, carbon credits \\
Logs and Alerts & 4 & Activity history, automated alert lifecycle \\
Irrigation & 3 & Schedule management, logs, AI recommendations \\
Finance & 11 & Costs, revenue, P\&L, market prices, schemes, insurance \\
Planning & 10 & Calendar, inventory, labour, equipment, rotation \\
Chat (LLM) & 2 & Gemini AI query, session history \\
Health Probes & 2 & Liveness + readiness for containers \\
\textbf{TOTAL} & \textbf{45} & \textbf{Full agricultural management coverage} \\
\hline
\end{tabularx}
\end{table}

\subsection{6.1.3. Database Design}
The database schema has 22 tables organized across 8 functional modules with full referential integrity in Third Normal Form (3NF).

Table 11: Core Entity-Relationship Model

\begin{table}[htbp]
\centering
\small
\begin{tabularx}{\textwidth}{|X|X|X|}
\hline
Entity & Key Fields & Relationships \\
\hline
User & id, email, username, password\_hash & 1:N with Fields, Costs, Revenue, ChatSessions \\
FieldData & id, user\_id, name, crop\_type, polygon & 1:N with HealthScans, Alerts, IrrigationLogs \\
CropHealthScan & id, field\_id, image\_path, disease, confidence & N:1 with FieldData \\
Cost & id, user\_id, field\_id, category, amount, date & N:1 with User, N:1 with FieldData \\
Revenue & id, user\_id, field\_id, source, amount, date & N:1 with User, N:1 with FieldData \\
FieldAlert & id, field\_id, message, is\_read, created\_at & N:1 with FieldData \\
FieldLog & id, user\_id, field\_id, activity, date, details & N:1 with User, N:1 with FieldData \\
IrrigationLog & id, field\_id, date, water\_amount, source & N:1 with FieldData \\
ChatSession & id, user\_id, created\_at, language\_code & N:1 with User \\
ChatMessage & id, session\_id, role, content, timestamp & N:1 with ChatSession \\
\hline
\end{tabularx}
\end{table}

\section{6.2. DESIGN NOTATIONS}
The following design notations are used throughout this chapter:

\begin{itemize}[leftmargin=*]
  \item \textbf{DFD (Data Flow Diagram):} Used to show how data moves through the system. Circles represent processes, rectangles represent external entities, open-ended rectangles represent data stores, and arrows represent data flows.
  \item \textbf{ER Diagram:} Used to show the database entity relationships. Rectangles represent entities, diamonds represent relationships, and lines with cardinality notation (1:N, N:1) show how entities relate.
  \item \textbf{Flowcharts:} Used to show the step-by-step logic of key algorithms. Rounded rectangles represent start/end, rectangles represent processes, diamonds represent decisions, and arrows represent flow direction.
  \item \textbf{ASCII Architecture Diagrams:} Used to show the system architecture, module dependencies, and deployment topology. Boxes represent components and arrows represent communication or dependency.
\end{itemize}
\section{6.3. DETAILED DESIGN}
\subsection{6.3.1. Machine Learning Pipeline Design}
\textbf{CNN Disease Detection Pipeline:}

The MobileNetV2 CNN uses ImageNet pre-trained weights with a custom 38-class output head. The architecture uses inverted residual blocks with depthwise separable convolutions, which makes it very lightweight (3.4M parameters, 9.1 MB).

Training configuration:

\begin{itemize}[leftmargin=*]
  \item Dataset: PlantVillage (87,848 augmented images, 80/20 train-test split)
  \item Optimizer: Adam (betas = 0.9, 0.999)
  \item Learning rate: 0.001, decayed by 0.5 every 5 epochs (StepLR)
  \item Batch size: 64; Total epochs: 10
  \item Progressive unfreezing: Layers frozen epochs 1-6; final conv block + classifier unfrozen epochs 7-10
  \item Augmentation: Random horizontal flip, rotation (15 degrees), color jitter, random resized crop
\end{itemize}
\textbf{LSTM Risk Prediction Pipeline:}

The two-layer LSTM has 64 hidden units per layer with 0.1 dropout and sigmoid output.

\begin{itemize}[leftmargin=*]
  \item Input: 10-day sequences of 4 features [NDVI, rainfall\_mm, temperature, soil\_moisture]
  \item Output: Binary risk probability p in [0, 1]
  \item Total parameters: 51,265 (about 208 KB)
  \item Training data: ERA5 climate reanalysis from Open-Meteo for 25 Indian districts (Jan 2023 - Dec 2024)
  \item 18,275 daily observations produce 18,025 sequences after 10-day windowing
  \item Class distribution: High-risk 54.6\%, Low-risk 45.4\% (naturally balanced)
\end{itemize}
\textbf{Composite Health Score:}

The health engine aggregates five weighted signals onto a common H in [0, 100] scale:

\begin{Verbatim}[fontsize=\small]
H = w1*S_CNN + w2*NDVI_norm + w3*(1 - p_LSTM)*100 + w4*S_weather + w5*S_practice
\end{Verbatim}
Where w1 = 0.30, w2 = 0.25, w3 = 0.20, w4 = 0.15, w5 = 0.10.

Rating thresholds: Excellent (>=80), Good (>=60), Fair (>=40), Poor (>=20), Critical (<20).

\subsection{6.3.2. API Design}
All 45 REST endpoints follow RESTful conventions with DRF class-based views:

\begin{itemize}[leftmargin=*]
  \item \textbf{Authentication:} Token-based with IsAuthenticated permission class
  \item \textbf{Serialization:} Explicit DRF serializers for all request/response bodies
  \item \textbf{Pagination:} PageNumberPagination (default page\_size=20) for list endpoints
  \item \textbf{Error Handling:} Centralized exception handler returning structured error format
\end{itemize}
\subsection{6.3.3. Security Design}
\begin{itemize}[leftmargin=*]
  \item RFC 6750 Token Authentication with 24-hour expiry
  \item PBKDF2-SHA256 hashing with 600,000 iterations
  \item Rate limiting: 5 login attempts/minute/IP
  \item Strict CORS origin whitelisting in production
  \item DRF serializer-based input validation
  \item File uploads restricted to image MIME types, 10 MB limit
  \item SQL injection prevention through exclusive Django ORM usage (zero raw SQL)
  \item PII encryption at rest using Fernet symmetric encryption
  \item HTTPS enforcement with TLS 1.3
\end{itemize}
\subsection{6.3.4. AWD Detection and Carbon Credit Design}
AWD cycle detection works by analyzing NDWI time-series for characteristic wet-dry-wet patterns:

\begin{itemize}[leftmargin=*]
  \item Wet threshold: NDWI > 0.3 (flooded/saturated field)
  \item Dry threshold: NDWI < 0.2 (drained/dry field)
  \item Minimum 1 complete cycle needed to confirm AWD practice
\end{itemize}
Carbon credit estimation uses IPCC Tier 2 methodology:

\begin{Verbatim}[fontsize=\small]
CH4_reduced = (EF_continuous - EF_AWD) * Area_ha * AWD_days
CO2e = CH4_reduced * GWP_100 / 1000
credit_value_usd = CO2e * carbon_price_per_tonne
\end{Verbatim}
Constants: EF\_continuous = 1.30 kg CH4/ha/day, EF\_AWD = 0.55 kg CH4/ha/day, GWP = 27.9, carbon price = USD 15/tonne.

\section{6.4. FLOWCHARTS}
\subsection{6.4.1. Disease Detection Workflow}
\begin{Verbatim}[fontsize=\small]
User uploads crop image via React frontend
         |
         v
+---------------------+
| Image Upload &      |
| MIME Type Validation |-- Invalid format --> Return HTTP 400
| (<=10 MB, image/*)  |
+---------------------+
         | Valid image
         v
+---------------------+
| Stage 1: Gemini     |
| Vision Pre-filter   |-- Non-plant image --> Return rejection message
| (Semantic Check)    |-- API unavailable --> Skip to Stage 2
+---------------------+
         | Valid plant image
         v
+---------------------+
| Stage 2: Image      |
| Preprocessing       |
| - Resize to 224x224 |
| - Normalize (ImageNet|
|   mean/std)         |
+---------------------+
         |
         v
+---------------------+
| MobileNetV2 CNN     |
| Forward Pass        |
| - 38-class softmax  |
| - Top-3 predictions |
+---------------------+
         |
         v
+---------------------+
| Output Response:    |
| - Disease class     |
| - Confidence score  |
| - Top-3 predictions |
| - Recommendations   |
+---------------------+
\end{Verbatim}
Figure 2: Disease Detection Workflow Flowchart

\subsection{6.4.2. Health Score Computation Workflow}
\begin{Verbatim}[fontsize=\small]
+----------+  +-----------+  +----------+  +----------+  +----------+
|  Crop    |  | Satellite |  | Weather  |  | LSTM     |  | NDWI/AWD |
|  Image   |  | NDVI Data |  | Data     |  | Risk     |  | Practice |
| (w1=0.30)|  | (w2=0.25) |  | (w4=0.15)|  | (w3=0.20)|  | (w5=0.10)|
+----+-----+  +-----+-----+  +----+-----+  +----+-----+  +----+-----+
     |              |             |             |             |
     v              v             v             v             v
  CNN Score    NDVI Norm    Weather Score  (1-risk)*100  Practice Score
     |              |             |             |             |
     +-------+------+------+-----+------+------+-------------+
                            |
                            v
                    +---------------+
                    | Weighted      |
                    | Fusion        |
                    | H in [0, 100] |
                    +-------+-------+
                            |
                            v
                    +---------------+
                    | Rating:       |
                    | Excellent     |
                    | Good / Fair   |
                    | Poor / Crit.  |
                    +---------------+
\end{Verbatim}
Figure 3: Health Score Computation Workflow

\subsection{6.4.3. User Authentication Flow}
\begin{Verbatim}[fontsize=\small]
Client                  Django Backend              PostgreSQL
  |                         |                          |
  |--- POST /signup ------->|                          |
  |    {username, email,    |--- CREATE user --------->|
  |     password}           |<-- User created ---------|
  |                         |--- Generate token ------>|
  |<-- 201 {token, user} ---|<-- Token stored ---------|
  |                         |                          |
  |--- POST /login -------->|                          |
  |    {username, password} |--- Verify credentials -->|
  |                         |<-- User found -----------|
  |                         |--- Check rate limit      |
  |<-- 200 {token, user} ---|                          |
  |                         |                          |
  |--- GET /field/data ---->|                          |
  |    Auth: Token xxx      |--- Validate token ------>|
  |                         |--- Query user fields --->|
  |<-- 200 {fields: [...]}--|<-- Results --------------|
\end{Verbatim}
Figure 4: User Authentication Flow

\section{6.5. PSEUDO CODE}
\subsection{6.5.1. CNN Prediction Algorithm}
\begin{Verbatim}[fontsize=\small]
FUNCTION predict_health(image_path):
    model = get_loaded_model()       // Lazy loading pattern
    IF model is None:
        RETURN fallback_response     // Model unavailable

    image = open_image(image_path)
    tensor = preprocess(image)       // Resize 224x224, normalize

    WITH no_gradient:
        output = model.forward(tensor)

    probabilities = softmax(output)
    top_class = argmax(probabilities)
    top_confidence = max(probabilities)

    class_name = CLASS_NAMES[top_class]
    is_healthy = top_class IN HEALTHY_INDICES
    healthy_prob = SUM(probabilities[i] for i in HEALTHY_INDICES)

    top3 = get_top_k(probabilities, k=3)

    RETURN {
        predicted_disease: class_name,
        probability: healthy_prob,
        confidence: categorize(top_confidence),   // High/Medium/Low
        top_predictions: top3
    }
END FUNCTION
\end{Verbatim}
\subsection{6.5.2. LSTM Risk Prediction Algorithm}
\begin{Verbatim}[fontsize=\small]
FUNCTION predict_risk(sequence_data):
    model, scaler = get_model_and_scaler()
    IF model is None OR scaler is None:
        RETURN fallback_risk         // Risk = 0.5, level = Unknown

    array = convert_to_numpy(sequence_data)
    array = replace_nan(array, 0.0)
    scaled = scaler.transform(array)
    tensor = to_tensor(scaled).unsqueeze(0)

    WITH no_gradient:
        probability = model.forward(tensor).item()

    IF probability > 0.7:
        level = "High"
        advice = "Immediate attention required"
    ELSE IF probability > 0.4:
        level = "Medium"
        advice = "Monitor closely"
    ELSE:
        level = "Low"
        advice = "Conditions favorable"

    RETURN {risk_probability: probability, risk_level: level, recommendation: advice}
END FUNCTION
\end{Verbatim}
\subsection{6.5.3. Carbon Credit Calculation Algorithm}
\begin{Verbatim}[fontsize=\small]
FUNCTION calculate_carbon_metrics(area_ha, awd_days, num_cycles):
    EF_CONTINUOUS = 1.30     // kg CH4/ha/day
    EF_AWD = 0.55            // kg CH4/ha/day
    GWP = 27.9               // IPCC AR6
    PRICE = 15.0             // USD/tonne CO2e
    WATER_SAVINGS = 525.0    // liters/ha/day

    ch4_reduced = (EF_CONTINUOUS - EF_AWD) * area_ha * awd_days
    co2e_tonnes = (ch4_reduced * GWP) / 1000
    credit_value = co2e_tonnes * PRICE
    water_saved = WATER_SAVINGS * area_ha * awd_days

    RETURN {ch4_reduced, co2e_tonnes, credit_value, water_saved}
END FUNCTION
\end{Verbatim}
\subsection{6.5.4. AWD Detection Algorithm}
\begin{Verbatim}[fontsize=\small]
FUNCTION detect_awd_from_ndwi(ndwi_series):
    WET_THRESHOLD = 0.3
    DRY_THRESHOLD = 0.2
    MIN_CYCLES = 1

    cycles = 0
    in_dry = FALSE

    FOR EACH value IN ndwi_series:
        IF value < DRY_THRESHOLD:
            in_dry = TRUE
        ELSE IF value > WET_THRESHOLD AND in_dry:
            cycles = cycles + 1
            in_dry = FALSE

    RETURN (cycles >= MIN_CYCLES, cycles)
END FUNCTION
\end{Verbatim}
\textbackslash\{\}newpage

\begin{center}
{\fontsize{14}{17}\bfseries 7. TESTING}\\
\end{center}
\section{7.1. FUNCTIONAL TESTING}
Functional testing verifies that each feature of the system works as expected from the user's perspective. We tested all 45 API endpoints and all 21 frontend dashboard modules against their specified requirements.

\subsection{7.1.1. API Endpoint Testing}
Every REST endpoint was tested with valid inputs, invalid inputs, missing authentication, and edge cases. The results are documented in the test results file (krishisaarthi\_results.txt).

Table 12: API Functional Test Results Summary

\begin{table}[htbp]
\centering
\small
\begin{tabularx}{\textwidth}{|X|X|X|X|X|}
\hline
Module & Endpoints Tested & Tests Passed & Tests Failed & Pass Rate \\
\hline
Authentication & 5 & 5 & 0 & 100\% \\
Field Management & 8 & 8 & 0 & 100\% \\
Logs and Alerts & 4 & 4 & 0 & 100\% \\
Irrigation & 3 & 3 & 0 & 100\% \\
Finance & 11 & 11 & 0 & 100\% \\
Planning & 10 & 10 & 0 & 100\% \\
Chat & 2 & 2 & 0 & 100\% \\
Health Probes & 2 & 2 & 0 & 100\% \\
\textbf{TOTAL} & \textbf{45} & \textbf{45} & \textbf{0} & \textbf{100\%} \\
\hline
\end{tabularx}
\end{table}

\subsection{7.1.2. ML Model Functional Testing}
\textbf{CNN Disease Detection:}

Table 13: CNN Functional Test Cases

\begin{table}[htbp]
\centering
\small
\begin{tabularx}{\textwidth}{|X|X|X|X|X|}
\hline
Test Case & Input & Expected Output & Actual Output & Result \\
\hline
Valid diseased leaf & Tomato Early Blight image & Disease class + confidence & Tomato Early Blight (0.96) & Pass \\
Valid healthy leaf & Healthy apple leaf & "Healthy" class & Healthy (0.99) & Pass \\
Non-plant image & Photo of a car & Rejected by Gemini Vision & "Not a plant image" & Pass \\
Corrupted file & 0-byte file & Error response & "Image file not found" & Pass \\
Oversized image & 15 MB photo & HTTP 400 & "File too large" & Pass \\
No auth token & Valid image, no token & HTTP 401 & "Authentication required" & Pass \\
\hline
\end{tabularx}
\end{table}

\textbf{LSTM Risk Prediction:}

Table 14: LSTM Functional Test Cases

\begin{table}[htbp]
\centering
\small
\begin{tabularx}{\textwidth}{|X|X|X|X|X|}
\hline
Test Case & Input & Expected Output & Actual Output & Result \\
\hline
High risk conditions & High temp, low moisture, low NDVI & Risk > 0.7, "High" & 0.87, "High" & Pass \\
Normal conditions & Moderate all values & Risk 0.3-0.5, "Medium" or "Low" & 0.38, "Low" & Pass \\
Missing data & NaN values in sequence & Handled gracefully & 0.50, fallback & Pass \\
Empty sequence & Empty list & Error response & "Empty sequence" & Pass \\
\hline
\end{tabularx}
\end{table}

\subsection{7.1.3. Frontend Functional Testing}
All 21 dashboard modules were manually tested for:

\begin{itemize}[leftmargin=*]
  \item Correct data display and formatting
  \item Form validation and error handling
  \item Navigation and routing
  \item Responsive layout on desktop and mobile viewports
  \item Language switching between all 4 supported languages
  \item Voice input functionality (on Chrome)
  \item Dark/light theme toggle
\end{itemize}
\section{7.2. STRUCTURAL TESTING}
Structural testing (white-box testing) examines the internal code structure to make sure all code paths are exercised.

\subsection{7.2.1. Code Coverage}
We used pytest with coverage tracking for the backend:

Table 15: Backend Code Coverage

\begin{table}[htbp]
\centering
\small
\begin{tabularx}{\textwidth}{|X|X|X|X|}
\hline
Module & Statements & Covered & Coverage \\
\hline
field/views & 456 & 389 & 85.3\% \\
finance/views & 312 & 271 & 86.9\% \\
planning/views & 287 & 248 & 86.4\% \\
chat/views & 98 & 82 & 83.7\% \\
ml\_engine/ & 234 & 198 & 84.6\% \\
KrishiSaarthi/ (settings, urls) & 89 & 78 & 87.6\% \\
\textbf{Total} & \textbf{1,476} & \textbf{1,266} & \textbf{85.8\%} \\
\hline
\end{tabularx}
\end{table}

\subsection{7.2.2. Path Testing}
Key decision paths tested include:

\begin{itemize}[leftmargin=*]
  \item CNN model loading: multiclass path, binary fallback path, no model path
  \item LSTM model loading: successful load, missing model file, missing scaler file
  \item Earth Engine API: successful call, API failure with circuit breaker, fallback NDVI
  \item Gemini Vision: plant image accepted, non-plant rejected, API unavailable (bypass)
  \item Authentication: valid token, expired token, missing token, rate limit exceeded
\end{itemize}
\subsection{7.2.3. Branch Coverage}
All conditional branches in the ML inference pipeline were tested:

\begin{itemize}[leftmargin=*]
  \item Model type checks (multiclass vs binary vs none)
  \item Confidence level thresholds (High > 0.8, Medium > 0.5, Low)
  \item Risk level thresholds (High > 0.7, Medium > 0.4, Low)
  \item Health score rating thresholds (Excellent >= 80, Good >= 60, Fair >= 40, Poor >= 20, Critical < 20)
\end{itemize}
\section{7.3. LEVELS OF TESTING}
\subsection{7.3.1. Unit Testing}
87 pytest test cases cover individual functions and methods:

\begin{itemize}[leftmargin=*]
  \item Model loading and initialization functions
  \item Data preprocessing and normalization
  \item Prediction output format validation
  \item Serializer validation for all API models
  \item Utility functions (coordinate validation, polygon validation)
  \item Carbon credit calculation with known inputs
\end{itemize}
\subsection{7.3.2. Integration Testing}
Integration tests verify that different modules work together correctly:

\begin{itemize}[leftmargin=*]
  \item User creates a field, then runs Earth Engine analysis on it
  \item User uploads an image, system runs Gemini validation then CNN inference
  \item Health score computation pulls data from CNN, LSTM, Earth Engine, and weather modules
  \item Financial entries are correctly aggregated into P\&L summaries
  \item Chat module injects field-specific context into Gemini queries
\end{itemize}
\subsection{7.3.3. System Testing}
End-to-end system tests simulate complete user workflows:

\begin{itemize}[leftmargin=*]
  \item Register, login, create field, run analysis, view dashboard, log out
  \item Upload disease image, get diagnosis, check health score, ask chatbot for advice
  \item Record costs and revenue, view P\&L, check market prices, discover government schemes
  \item Plan crop rotation, manage inventory, schedule equipment, track labor
\end{itemize}
\subsection{7.3.4. Performance Testing}
API endpoints were benchmarked under realistic conditions:

Table 16: API Endpoint Latency Profile

\begin{table}[htbp]
\centering
\small
\begin{tabularx}{\textwidth}{|X|X|X|X|}
\hline
Endpoint & Method & Mean Latency & Status \\
\hline
POST /login & POST & 45 ms & Pass \\
GET /field/list & GET & 12 ms & Pass \\
POST /field/diagnose & POST & 1,850 ms & Pass (includes Gemini Vision) \\
GET /field/health-score & GET & 234 ms & Pass (includes CNN + LSTM) \\
GET /field/ee-analysis & GET & 3,200 ms & Pass (Earth Engine API latency) \\
POST /chat/query & POST & 2,100 ms & Pass (Gemini Chat latency) \\
GET /finance/pnl & GET & 18 ms & Pass \\
GET /finance/market-prices & GET & 890 ms & Pass (external API) \\
GET /field/carbon-credits & GET & 156 ms & Pass \\
GET /planning/calendar & GET & 15 ms & Pass \\
GET /health/live & GET & 3 ms & Pass \\
GET /health/ready & GET & 45 ms & Pass (ML model check) \\
\hline
\end{tabularx}
\end{table}

\section{7.4. TESTING THE PROJECT}
\subsection{7.4.1. ML Model Performance Results}
\textbf{CNN Crop Disease Classifier:}

Table 17: CNN Overall Performance Metrics

\begin{table}[htbp]
\centering
\small
\begin{tabularx}{\textwidth}{|X|X|}
\hline
Metric & Value \\
\hline
Top-1 Accuracy & 99.03\% \\
Weighted Precision & 0.9905 \\
Weighted Recall & 0.9903 \\
Weighted F1-Score & 0.9903 \\
Test Loss & 0.0411 \\
\hline
\end{tabularx}
\end{table}

Table 18: CNN Training Convergence Per-Epoch

\begin{table}[htbp]
\centering
\small
\begin{tabularx}{\textwidth}{|X|X|X|X|X|X|}
\hline
Epoch & Train Loss & Train Acc (\%) & Test Loss & Test Acc (\%) & Phase \\
\hline
1 & 0.4521 & 87.23 & 0.2103 & 93.45 & Frozen \\
2 & 0.1876 & 94.12 & 0.1245 & 96.18 & Frozen \\
3 & 0.1204 & 96.34 & 0.0891 & 97.42 & Frozen \\
4 & 0.0934 & 97.11 & 0.0723 & 97.89 & Frozen \\
5 & 0.0781 & 97.56 & 0.0645 & 98.12 & Frozen \\
6 & 0.0689 & 97.89 & 0.0601 & 98.34 & Frozen \\
7 & 0.0598 & 98.12 & 0.0534 & 98.56 & Unfrozen \\
8 & 0.0467 & 98.45 & 0.0489 & 98.78 & Unfrozen \\
9 & 0.0389 & 98.71 & 0.0445 & 98.91 & Unfrozen \\
10 & 0.0312 & 98.89 & 0.0411 & 99.03 & Unfrozen \\
\hline
\end{tabularx}
\end{table}

The progressive unfreezing strategy clearly shows its value. When the final convolutional block is unfrozen at epoch 7, accuracy jumps from 98.34\% to 99.03\% over the final 4 epochs.

Table 19: CNN Inference Performance

\begin{table}[htbp]
\centering
\small
\begin{tabularx}{\textwidth}{|X|X|}
\hline
Metric & Value \\
\hline
Mean Inference Latency (CPU) & 0.17 ms \\
Model Size (FP32) & 9.1 MB \\
Parameters & 3.4 million \\
Memory Footprint & \textasciitilde{}45 MB \\
Throughput & \textasciitilde{}5,800 images/second \\
\hline
\end{tabularx}
\end{table}

\textbf{LSTM Risk Forecasting Model:}

Table 20: LSTM Risk Model Performance

\begin{table}[htbp]
\centering
\small
\begin{tabularx}{\textwidth}{|X|X|}
\hline
Metric & Value \\
\hline
Accuracy & 97.56\% \\
Precision & 0.9785 \\
Recall & 0.9724 \\
F1-Score (weighted) & 0.9776 \\
AUC-ROC & 0.9983 \\
Test Loss (BCE) & 0.0712 \\
\hline
\end{tabularx}
\end{table}

\subsection{7.4.2. Composite Health Score Validation}
Table 21: Health Score Validation Scenarios

\begin{table}[htbp]
\centering
\small
\begin{tabularx}{\textwidth}{|X|X|X|X|X|X|X|X|}
\hline
Scenario & CNN & NDVI & LSTM Risk & Weather & Practice & Health Score & Rating \\
\hline
Best Case & 95 & 100 & 0.05 & 90 & 100 & 95.0 & Excellent \\
Good & 80 & 75 & 0.15 & 75 & 60 & 74.5 & Good \\
Mixed & 50 & 60 & 0.40 & 60 & 40 & 55.0 & Fair \\
Stressed & 30 & 35 & 0.70 & 40 & 20 & 34.3 & Poor \\
Critical & 10 & 15 & 0.95 & 15 & 0 & 12.0 & Critical \\
\hline
\end{tabularx}
\end{table}

The 83-point range (95.0 to 12.0) across scenarios confirms the system's ability to tell apart different agricultural conditions effectively.

\subsection{7.4.3. AWD and Carbon Credit Validation}
Table 22: AWD Detection Accuracy

\begin{table}[htbp]
\centering
\small
\begin{tabularx}{\textwidth}{|X|X|X|X|X|}
\hline
Scenario & Field Size & AWD Cycles & AWD Days & Detection \\
\hline
Active AWD (irrigated rice) & 2.0 ha & 3 & 45 & Correct \\
Continuous flooding & 2.0 ha & 0 & 0 & Correct \\
Partial AWD (1 cycle) & 1.5 ha & 1 & 15 & Correct \\
Rainfed (no irrigation) & 1.0 ha & 0 & 0 & Correct \\
\hline
\end{tabularx}
\end{table}

Table 23: Carbon Credit Estimation for 2-ha Field with 3 AWD Cycles (45 days)

\begin{table}[htbp]
\centering
\small
\begin{tabularx}{\textwidth}{|X|X|}
\hline
Metric & Value \\
\hline
CH4 Reduction & 67.50 kg \\
CO2-equivalent & 1.88 tonnes \\
Credit Value (USD) & \$28.22 \\
Credit Value (INR) & Rs 2,369 \\
Water Saved & 47,250 liters \\
\hline
\end{tabularx}
\end{table}

\subsection{7.4.4. System Resource Utilization}
Table 24: System Resource Utilization

\begin{table}[htbp]
\centering
\small
\begin{tabularx}{\textwidth}{|X|X|X|}
\hline
Resource & Idle & Under Load (10 concurrent requests) \\
\hline
CPU Usage & 2-5\% & 35-45\% \\
Memory (Backend) & 180 MB & 420 MB (with models loaded) \\
Memory (Frontend) & - & 85 MB (static serving) \\
Memory (PostgreSQL) & 50 MB & 120 MB \\
Disk (Total Deployment) & 2.1 GB & - \\
Disk (ML Models) & 9.3 MB & - \\
\hline
\end{tabularx}
\end{table}

\textbackslash\{\}newpage

\begin{center}
{\fontsize{14}{17}\bfseries 8. IMPLEMENTATION}\\
\end{center}
\section{8.1. IMPLEMENTATION OF THE PROJECT}
\subsection{8.1.1. Backend Implementation}
The backend was implemented using Django 5.1 with Django REST Framework. The codebase is organized into 8 Django apps, each responsible for a specific functional domain:

\textbf{Authentication Module:} Handles user registration, login, logout, and password reset. Uses Django's built-in User model with TokenAuthentication from DRF. Rate limiting is implemented using a custom middleware that tracks login attempts per IP address.

\textbf{Field Module:} The core module that manages field data, satellite analysis, disease detection, health scores, AWD detection, and carbon credit estimation. Earth Engine integration uses the \texttt{ee} Python library with a service account for authentication. The circuit breaker pattern protects against API failures.

\textbf{Finance Module:} Manages costs, revenue, seasons, P\&L calculations, market prices, government schemes, and insurance claims. Market prices are fetched from external APIs and cached using Redis (with LocMemCache fallback). P\&L calculations use database-level aggregation (Django's Sum and Count) instead of Python-side loops for efficiency.

\textbf{Planning Module:} Handles crop calendar, inventory tracking, labor management, equipment scheduling, and crop rotation planning. Each module uses DRF serializers for input validation and ModelViewSets for standard CRUD operations.

\textbf{Chat Module:} Integrates Google Gemini 2.5 Flash for conversational AI. Each query is enriched with field-specific context (NDVI, weather, risk level) before being sent to Gemini. Conversation history is maintained per session.

\textbf{ML Engine:} A standalone Python module containing the CNN, LSTM, health score, AWD detection, and carbon credit calculation logic. Models are lazy-loaded on first use to reduce startup time.

\subsection{8.1.2. Frontend Implementation}
The frontend was built using React 18 with TypeScript and Vite as the build tool. Key implementation decisions:

\textbf{Component Architecture:} 79 components organized across 21 dashboard modules. Each module has its own directory with component files, types, and styles. Shadcn UI provides the base component library, customized with TailwindCSS 4 design tokens.

\textbf{State Management:} React Context API is used for global state (authentication, selected field, theme). Local component state uses useState and useReducer hooks. No external state management library (like Redux) was needed because the application state is relatively straightforward.

\textbf{API Layer:} A centralized api.ts utility handles all backend communication with automatic token injection, retry logic with exponential backoff, and error handling. All API calls go through this layer, ensuring consistent behavior.

\textbf{Maps:} Leaflet.js with react-leaflet provides interactive maps for field boundary drawing. Users can draw polygons on satellite imagery to define their field boundaries. The coordinates are stored as GeoJSON in the backend.

\textbf{Internationalization:} A custom useTranslation hook provides multilingual support. Translation strings are stored in JSON files for each language. Language switching happens without page reload.

\textbf{Voice Input:} A useVoice hook wraps the Web Speech API SpeechRecognition interface. Voice input is available in the chat module for hands-free interaction.

\subsection{8.1.3. ML Model Implementation}
\textbf{CNN Training:} The MobileNetV2 model was trained on the PlantVillage dataset using the progressive unfreezing strategy described in Chapter 6. Training was done entirely on CPU (Intel i7, no GPU), taking approximately 4 hours for 10 epochs with the 87,848-image dataset.

\textbf{LSTM Training:} The risk prediction model was trained on ERA5 climate data collected for 25 Indian districts. Data collection used the Open-Meteo Historical Weather API. The model converged within 5 epochs but was trained for 30 epochs to verify stability.

\textbf{Model Registry:} All trained models are tracked via a registry module that provides versioning, SHA-256 integrity verification at startup, and health status reporting through the readiness endpoint.

\section{8.2. CONVERSION PLAN}
\subsection{8.2.1. Development to Production Conversion}
The conversion from development to production involves:

\begin{enumerate}[leftmargin=*]
  \item \textbf{Database Migration:} Switch from SQLite to PostgreSQL using Django's database router. Run \texttt{python manage.py migrate} against the production database. Initial data seeding via the init\_db.py script.
  \item \textbf{Static Files:} Run \texttt{python manage.py collectstatic} to gather all static files. Nginx serves static files directly, bypassing Django for better performance.
  \item \textbf{WSGI Server:} Replace Django's development server with Gunicorn (4 workers, 2 threads per worker). This provides proper process management, graceful restarts, and request queuing.
  \item \textbf{Reverse Proxy:} Nginx handles TLS termination, static file serving, request routing, and gzip compression. The nginx.conf file is configured for production security headers.
  \item \textbf{Environment Variables:} All secrets (Django secret key, API keys) are loaded from environment variables via python-dotenv. The .env.example file documents all required variables.
  \item \textbf{Docker Deployment:} Docker Compose orchestrates all services (backend, frontend, database, Redis, Prometheus, Grafana). Production uses docker-compose.prod.yml with optimized configurations.
\end{enumerate}
\subsection{8.2.2. Container Architecture}
Table 25: Container Architecture

\begin{table}[htbp]
\centering
\small
\begin{tabularx}{\textwidth}{|X|X|X|X|}
\hline
Service & Image Base & Port & Purpose \\
\hline
backend & python:3.11-slim & 8000 & Django API (Gunicorn) \\
frontend & node:20-alpine then nginx:alpine & 80/443 & React SPA \\
db & postgres:14-alpine & 5432 & PostgreSQL database \\
redis & redis:7-alpine & 6379 & Session cache \\
prometheus & prom/prometheus & 9090 & Metrics collection \\
grafana & grafana/grafana & 3000 & Metrics visualization \\
\hline
\end{tabularx}
\end{table}

\section{8.3. POST-IMPLEMENTATION AND SOFTWARE MAINTENANCE}
\subsection{8.3.1. Monitoring and Observability}
The platform includes a complete observability stack:

\begin{itemize}[leftmargin=*]
  \item \textbf{Prometheus} collects metrics including request latency histograms, error rate counters, and ML inference duration.
  \item \textbf{Grafana} dashboards visualize API latency, error rates, database performance, and system resource usage.
  \item \textbf{Health Probes:} /health/live (liveness) and /health/ready (readiness with ML model checks) support container orchestration platforms like Kubernetes.
  \item \textbf{Structured Logging:} JSON-formatted logs with correlation IDs enable request tracing across the entire pipeline.
\end{itemize}
\subsection{8.3.2. CI/CD Pipeline}
GitHub Actions automates the following on every push:

\begin{enumerate}[leftmargin=*]
  \item \textbf{Lint and Static Analysis:} Ruff for Python, ESLint/Prettier for TypeScript
  \item \textbf{Unit Tests:} pytest (backend) + Vitest (frontend)
  \item \textbf{Docker Build:} Multi-stage builds with layer caching
  \item \textbf{Integration Tests:} API endpoint validation
  \item \textbf{Security Scan:} pip-audit (Python) + npm audit (Node.js)
  \item \textbf{Deploy:} Docker Compose deployment with zero-downtime restart
\end{enumerate}
\subsection{8.3.3. Maintenance Plan}
\begin{itemize}[leftmargin=*]
  \item \textbf{Model Retraining:} CNN model should be retrained when new crop species are added to the PlantVillage dataset or when field-captured images become available. LSTM model should be retrained annually with updated ERA5 data.
  \item \textbf{Dependency Updates:} Monthly review of Python and Node.js dependencies for security patches.
  \item \textbf{Database Maintenance:} Weekly backup schedule. Index optimization quarterly.
  \item \textbf{API Key Rotation:} Gemini API key and Earth Engine credentials should be rotated every 6 months.
  \item \textbf{User Feedback:} Structured feedback collection from pilot users to prioritize feature improvements.
\end{itemize}
\textbackslash\{\}newpage

\begin{center}
{\fontsize{14}{17}\bfseries 9. PROJECT LEGACY}\\
\end{center}
\section{9.1. CURRENT STATUS OF THE PROJECT}
As of April 2026, KrishiSaarthi is a fully functional, production-ready platform. Here is the current status of each component:

Table 26: Current Project Status

\begin{table}[htbp]
\centering
\small
\begin{tabularx}{\textwidth}{|X|X|X|}
\hline
Component & Status & Details \\
\hline
Backend API & Complete & 45 REST endpoints, all functional and tested \\
Frontend Dashboard & Complete & 79 components across 21 dashboard modules \\
CNN Disease Model & Complete & 99.03\% accuracy on 38-class PlantVillage test set \\
LSTM Risk Model & Complete & 97.56\% accuracy, AUC-ROC 0.9983 \\
Health Score Engine & Complete & 5-signal fusion, 83-point discriminative range \\
AWD/Carbon Credit & Complete & IPCC Tier 2 methodology, validated on test scenarios \\
Earth Engine Integration & Complete & 4 vegetation indices with circuit breaker fallback \\
Chat Advisory & Complete & Gemini 2.5 Flash, 4 languages, voice input \\
Financial Management & Complete & Full P\&L suite with market prices and schemes \\
Resource Planning & Complete & Calendar, inventory, labor, equipment, rotation \\
Docker Deployment & Complete & Dev and production Docker Compose configurations \\
CI/CD Pipeline & Complete & GitHub Actions with lint, test, build, security scan \\
Monitoring & Complete & Prometheus + Grafana stack configured \\
API Documentation & Complete & OpenAPI 3.0, Swagger UI, ReDoc \\
Source Code & Released & MIT License on GitHub \\
\hline
\end{tabularx}
\end{table}

The platform is deployed and running in a development/demonstration environment. It has not yet been deployed for real-world use by actual farmers. That is identified as the primary next step.

\section{9.2. REMAINING AREAS OF CONCERN}
\subsection{9.2.1. Field Validation Gap}
The most significant remaining concern is the lack of real-world field validation. The CNN model has been evaluated only on the PlantVillage laboratory dataset. Real-world crop images taken by farmers in uncontrolled lighting, angles, and backgrounds may produce lower accuracy. The two-stage Gemini Vision pre-filtering helps mitigate this, but proper field testing is needed.

\subsection{9.2.2. Limited Crop Coverage}
The current CNN model covers only 14 crop species (38 classes) from PlantVillage. Major Indian crops like rice, wheat, cotton, and sugarcane are not yet covered. Expanding coverage requires collecting and labeling field-captured images for these crops.

\subsection{9.2.3. AWD Ground Truth}
The AWD detection algorithm has been validated on synthetic and semi-real NDWI time-series. Comprehensive validation against ground-truth IoT sensor data has not been conducted. The carbon credit estimates should be treated as indicative rather than auditable until field validation is completed.

\subsection{9.2.4. Scalability Under Production Load}
The current performance benchmarks are based on a development server setup with SQLite and a single Django process. Production deployment with PostgreSQL, Gunicorn (4 workers), and Redis caching will improve performance, but load testing with hundreds of concurrent users has not been done.

\subsection{9.2.5. Internet Dependency}
The platform relies on internet connectivity for satellite data (Earth Engine), weather updates (OpenWeather), and AI advisory (Gemini). While fallback mechanisms exist for each external service, the platform cannot function in a fully offline mode. A mobile app with offline-first architecture is planned as future work.

\section{9.3. TECHNICAL AND MANAGERIAL LESSONS LEARNT}
\subsection{9.3.1. Technical Lessons}
\textbf{1. Transfer learning is extremely powerful for resource-constrained environments.} We achieved 99.03\% accuracy with just 10 epochs of training on CPU by leveraging ImageNet pre-trained weights and progressive unfreezing. This completely eliminated the need for GPU infrastructure, which would have added significant cost and complexity.

\textbf{2. Circuit breaker patterns are essential for external API reliability.} Without the circuit breaker, Earth Engine API outages (which happen roughly 2-3 times per month) would cascade into complete health score failures. The pattern allows graceful degradation with deterministic fallback values.

\textbf{3. Feature selection matters more than model complexity for time-series prediction.} The LSTM model achieves 97.56\% accuracy with just 4 input features and 51,265 parameters. Adding more features or increasing model complexity was not necessary because the selected features (NDVI, rainfall, temperature, soil moisture) already capture the essential risk signals.

\textbf{4. LLM context injection dramatically improves response relevance.} Simply asking Gemini "how to care for crops" gives generic answers. But injecting field-specific data (current NDVI of 0.45, last rainfall 3 days ago, current temperature 34 degrees C) into the system prompt produces highly specific and actionable advice.

\textbf{5. Modular Django architecture enables effective parallel development.} By dividing the backend into 8 independent Django apps (field, finance, planning, chat, etc.), team members could work on different modules simultaneously with minimal merge conflicts. However, cross-module data flows (like health score computation that pulls from multiple modules) required careful coordination.

\subsection{9.3.2. Managerial Lessons}
\textbf{1. Start with the ML pipeline, not the UI.} We initially spent too much time on frontend design before the ML models were trained. In hindsight, training the CNN and LSTM models first would have given us concrete outputs to design the UI around, rather than using placeholder data.

\textbf{2. External API dependencies introduce schedule risk.} Earth Engine authentication issues and Gemini API rate limits caused delays that were hard to predict. Building fallback mechanisms early (rather than as an afterthought) would have reduced stress during integration testing.

\textbf{3. Documentation during development saves time later.} We adopted OpenAPI auto-generation (drf-spectacular) early, which meant API documentation stayed current without manual effort. For code comments and architecture docs, we were not as disciplined, and catching up during the report writing phase was time-consuming.

\textbf{4. Regular milestone reviews keep the team focused.} Our bi-weekly milestone reviews helped identify delays early (e.g., M3 was 5 days late, M7 was 3 days late). Without these checkpoints, the delays could have compounded.

\textbf{5. Seven-member teams need clear module ownership.} With 7 team members, assigning clear ownership of specific modules (2-3 people per major module) worked better than having everyone touch everything. This reduced communication overhead and improved accountability.

\textbackslash\{\}newpage

\begin{center}
{\fontsize{14}{17}\bfseries 10. USER MANUAL}\\
\end{center}
\section{10.1. GETTING STARTED}
\subsection{10.1.1. System Requirements}
To use KrishiSaarthi, you need:

\begin{itemize}[leftmargin=*]
  \item A device (computer, tablet, or smartphone) with a web browser
  \item Internet connection
  \item A modern web browser: Chrome 90+, Firefox 90+, Safari 15+, or Edge 90+
\end{itemize}
\subsection{10.1.2. Accessing the Platform}
Open your web browser and go to the KrishiSaarthi URL provided by your administrator. For local development, this is typically http://localhost:5173.

\subsection{10.1.3. Creating an Account}
\begin{enumerate}[leftmargin=*]
  \item Click the "Sign Up" button on the landing page.
  \item Enter your desired username, email address, and password.
  \item Click "Create Account."
  \item You will be automatically logged in and taken to your dashboard.
\end{enumerate}
\subsection{10.1.4. Logging In}
\begin{enumerate}[leftmargin=*]
  \item Enter your username and password on the login page.
  \item Click "Login."
  \item If you have forgotten your password, click "Forgot Password" and follow the email instructions.
\end{enumerate}
\section{10.2. MANAGING YOUR FIELDS}
\subsection{10.2.1. Adding a New Field}
\begin{enumerate}[leftmargin=*]
  \item From the dashboard, click "Add Field" or go to the Field Management section.
  \item Enter a name for your field (e.g., "North Paddy Field").
  \item Select your crop type from the dropdown menu.
  \item Use the interactive map to draw the boundary of your field by clicking on the corners.
  \item Click "Save Field."
\end{enumerate}
\subsection{10.2.2. Viewing Field Details}
Click on any field card on your dashboard to see:

\begin{itemize}[leftmargin=*]
  \item Field name, crop type, and area
  \item Current vegetation indices (NDVI, EVI, SAVI, NDWI)
  \item Health score and rating
  \item Recent alerts and activity logs
\end{itemize}
\subsection{10.2.3. Running Satellite Analysis}
\begin{enumerate}[leftmargin=*]
  \item Select a field from your dashboard.
  \item Click "Run Analysis" or go to the Analytics section.
  \item The system will fetch the latest Sentinel-2 satellite data and compute vegetation indices.
  \item Results are displayed as charts showing trends over time.
\end{enumerate}
\section{10.3. DISEASE DETECTION}
\begin{enumerate}[leftmargin=*]
  \item Go to the "Diagnosis" section from the sidebar.
  \item Click "Upload Image" and select a clear photo of a crop leaf.
  \item The system will first verify that the image shows a plant leaf.
  \item If verified, the CNN model will analyze the image and return:
\begin{itemize}[leftmargin=*]
  \item The most likely disease (or "Healthy" if no disease is detected)
  \item Confidence level (High, Medium, or Low)
  \item Top 3 most likely diagnoses
  \item Recommended actions
\end{itemize}
\section{10.4. FINANCIAL MANAGEMENT}
\subsection{10.4.1. Recording Costs}
\begin{enumerate}[leftmargin=*]
  \item Go to "Finance" then "Costs" from the sidebar.
  \item Click "Add Cost."
  \item Select the category (seeds, fertilizer, labor, equipment, other).
  \item Enter the amount, date, and a brief description.
  \item Click "Save."
\end{enumerate}
\subsection{10.4.2. Recording Revenue}
\begin{enumerate}[leftmargin=*]
  \item Go to "Finance" then "Revenue."
  \item Click "Add Revenue."
  \item Enter the source, amount, date, and description.
  \item Click "Save."
\end{enumerate}
\subsection{10.4.3. Viewing Profit and Loss}
Go to "Finance" then "P\&L Dashboard" to see a visual summary of your costs, revenue, and net profit or loss for each season.

\subsection{10.4.4. Checking Market Prices}
Go to "Finance" then "Market Prices" to see current commodity prices from major markets.

\section{10.5. AI CHAT ADVISORY}
\begin{enumerate}[leftmargin=*]
  \item Click the chat icon in the sidebar or bottom-right corner.
  \item Type your question in your preferred language (English, Hindi, Punjabi, or Malayalam).
  \item Alternatively, click the microphone icon and speak your question.
  \item The AI will respond with personalized advice based on your field's current conditions.
\end{enumerate}
Example questions:

\begin{itemize}[leftmargin=*]
  \item "My tomato leaves have yellow spots. What should I do?"
  \item "When is the best time to apply fertilizer for wheat?"
  \item "What government schemes are available for rice farmers in Punjab?"
\end{itemize}
\section{10.6. CARBON CREDITS}
\begin{enumerate}[leftmargin=*]
  \item Go to the "Carbon Credits" section from the field details page.
  \item The system will analyze your irrigation patterns over the past 90 days.
  \item If AWD (Alternate Wetting and Drying) irrigation is detected, the system will show:
\begin{itemize}[leftmargin=*]
  \item Number of AWD cycles detected
  \item Estimated methane reduction (kg CH4)
  \item Carbon credit value in both USD and INR
  \item Water savings (liters)
\end{itemize}
\section{10.7. LANGUAGE AND SETTINGS}
\begin{itemize}[leftmargin=*]
  \item Click the language selector in the top bar to switch between English, Hindi, Punjabi, and Malayalam.
  \item Click the theme toggle to switch between light and dark modes.
  \item All settings are saved to your account and persist across sessions.
\end{itemize}
\textbackslash\{\}newpage

\begin{center}
{\fontsize{14}{17}\bfseries 11. SOURCE CODE AND SYSTEM SNAPSHOTS}\\
\end{center}
\section{11.1. SOURCE CODE HIGHLIGHTS}
\subsection{11.1.1. CNN Prediction Module (backend/ml_engine/cnn.py)}
\begin{Verbatim}[fontsize=\small]
# Image preprocessing pipeline
transform = transforms.Compose([
    transforms.Resize((224, 224)),
    transforms.ToTensor(),
    transforms.Normalize([0.485, 0.456, 0.406], [0.229, 0.224, 0.225]),
])

def predict_health(img_path, device="cpu"):
    """Predict crop health/disease from an image."""
    model = get_model()
    if model is None:
        return {"error": "Model not available", "probability": 0.5}

    img = Image.open(img_path).convert("RGB")
    img_tensor = transform(img).unsqueeze(0).to(device)

    with torch.no_grad():
        output = model(img_tensor)

    probs = torch.softmax(output, dim=1).squeeze(0)
    top_prob, top_idx = probs.max(0)
    predicted_class = CLASS_NAMES[top_idx.item()]
    is_healthy = top_idx.item() in HEALTHY_INDICES
    healthy_prob = sum(probs[i].item() for i in HEALTHY_INDICES)
    top3_probs, top3_indices = probs.topk(3)

    return {
        "probability": round(healthy_prob, 4),
        "class": "Healthy" if is_healthy else "Infested",
        "predicted_disease": predicted_class,
        "confidence_score": round(top_prob.item(), 4),
        "top_predictions": [
            {"class": CLASS_NAMES[idx.item()],
             "probability": round(p.item(), 4)}
            for p, idx in zip(top3_probs, top3_indices)
        ],
    }
\end{Verbatim}
\subsection{11.1.2. LSTM Risk Prediction Module (backend/ml_engine/lstm.py)}
\begin{Verbatim}[fontsize=\small]
class RiskLSTM(nn.Module):
    """LSTM model for predicting crop disease/pest risk."""

    def __init__(self, input_size=4, hidden_size=64,
                 num_layers=2, dropout=0.1):
        super().__init__()
        self.lstm = nn.LSTM(
            input_size, hidden_size, num_layers,
            batch_first=True,
            dropout=dropout if num_layers > 1 else 0,
        )
        self.fc = nn.Linear(hidden_size, 1)
        self.act = nn.Sigmoid()

    def forward(self, x):
        out, _ = self.lstm(x)
        out = out[:, -1, :]   # Take last timestep
        out = self.fc(out)
        return self.act(out)
\end{Verbatim}
\subsection{11.1.3. Health Score Computation (backend/ml_engine/health_score.py)}
\begin{Verbatim}[fontsize=\small]
WEIGHT_CNN = 0.30
WEIGHT_NDVI = 0.25
WEIGHT_LSTM = 0.20
WEIGHT_WEATHER = 0.15
WEIGHT_PRACTICE = 0.10

def compute_health_score(cnn_confidence, ndvi_value, lstm_risk,
                         weather_score, practice_score):
    """Five-signal weighted fusion onto H in [0, 100] scale."""
    cnn_score = (1 - cnn_confidence) * 100
    ndvi_score = min(ndvi_value / 0.8, 1.0) * 100
    lstm_score = (1 - lstm_risk) * 100

    health = (WEIGHT_CNN * cnn_score +
              WEIGHT_NDVI * ndvi_score +
              WEIGHT_LSTM * lstm_score +
              WEIGHT_WEATHER * weather_score +
              WEIGHT_PRACTICE * practice_score)

    return max(0, min(100, health))
\end{Verbatim}
\subsection{11.1.4. Carbon Credit Calculation (backend/ml_engine/cc.py)}
\begin{Verbatim}[fontsize=\small]
EF_CONTINUOUS = 1.30    # kg CH4/ha/day
EF_AWD = 0.55           # kg CH4/ha/day
GWP_CH4 = 27.9          # IPCC AR6 100-year GWP
CARBON_PRICE = 15.0     # USD/tonne CO2e
WATER_SAVINGS = 525.0   # liters/ha/day

def calculate_carbon_metrics(area_ha, awd_days, num_cycles):
    ch4_reduced = (EF_CONTINUOUS - EF_AWD) * area_ha * awd_days
    co2e_tonnes = (ch4_reduced * GWP_CH4) / 1000
    credit_value = co2e_tonnes * CARBON_PRICE
    water_saved = WATER_SAVINGS * area_ha * awd_days
    return {
        'ch4_reduced_kg': ch4_reduced,
        'co2e_tonnes': co2e_tonnes,
        'credit_value_usd': credit_value,
        'water_saved_liters': water_saved
    }
\end{Verbatim}
\subsection{11.1.5. Field Data Model (backend/field/models.py)}
\begin{Verbatim}[fontsize=\small]
class FieldData(models.Model):
    user = models.ForeignKey(User, on_delete=models.CASCADE,
                             related_name="fields")
    name = models.CharField(max_length=100, default="My Field")
    cropType = models.CharField(max_length=32)
    polygon = models.JSONField(validators=[validate_polygon])
    created_at = models.DateTimeField(auto_now_add=True)
    updated_at = models.DateTimeField(auto_now=True)

    class Meta:
        indexes = [models.Index(fields=["user", "created_at"])]

class CropHealthScan(models.Model):
    id = models.UUIDField(primary_key=True, default=uuid.uuid4)
    user = models.ForeignKey(User, on_delete=models.CASCADE,
                             related_name="health_scans")
    field = models.ForeignKey(FieldData, on_delete=models.CASCADE,
                              related_name="health_scans")
    image_path = models.ImageField(upload_to="ml_scans/")
    detected_disease = models.CharField(max_length=100, null=True)
    confidence_score = models.FloatField(null=True)
    recommendation = models.TextField(null=True)
    severity = models.CharField(max_length=20, default="LOW")
    scanned_at = models.DateTimeField(auto_now_add=True, db_index=True)
\end{Verbatim}
\section{11.2. PROJECT STRUCTURE}
\begin{Verbatim}[fontsize=\small]
KrishiSaarthi/
+-- backend/                      # Django REST API Backend
|   +-- KrishiSaarthi/           # Django project settings
|   |   +-- settings.py          # Configuration
|   |   +-- urls.py              # Root URL routing
|   +-- field/                   # Field management app
|   |   +-- models.py            # FieldData, Pest, Alert models
|   |   +-- views/               # API view classes
|   |   +-- serializers.py       # API serializers
|   |   +-- utils.py             # Earth Engine utilities
|   +-- finance/                 # Financial management app
|   +-- planning/                # Resource planning app
|   +-- chat/                    # AI assistant app
|   +-- ml_engine/               # Machine learning module
|   |   +-- cnn.py               # MobileNetV2 pest detection
|   |   +-- lstm.py              # Risk prediction model
|   |   +-- health_score.py      # Multi-factor health scoring
|   |   +-- awd.py               # AWD detection from NDWI
|   |   +-- cc.py                # Carbon credit calculations
|   |   +-- registry.py          # Model versioning and integrity
|   +-- ml_models/               # Trained model weights
|       +-- crop_health_model.pth   # 9.1MB CNN model
|       +-- risk_lstm_final.pth     # LSTM weights
|       +-- risk_scaler.save        # Input normalization
|
+-- frontend/                    # React TypeScript Frontend
|   +-- client/
|       +-- src/
|           +-- components/      # UI Components (79 files)
|           +-- pages/           # Dashboard, Landing, Auth
|           +-- context/         # Auth, Field, Theme providers
|           +-- hooks/           # useTranslation, useVoice
|           +-- lib/             # API client, utilities
|           +-- types/           # TypeScript interfaces
|
+-- docker-compose.yml           # Development orchestration
+-- docker-compose.prod.yml      # Production orchestration
+-- nginx/                       # Nginx configuration
+-- monitoring/                  # Prometheus/Grafana configs
\end{Verbatim}
\textbackslash\{\}newpage

\begin{center}
{\fontsize{14}{17}\bfseries 12. BIBLIOGRAPHY}\\
\end{center}
[1] Manzoor, T. (2024). "Precision Agriculture: A Comprehensive Review of Machine Learning Techniques." *Journal of Agricultural Informatics*, 15(2), 45-67.

[2] Sudha, S. V., \& Loret, P. (2026). "Integration of Machine Learning, IoT, and Blockchain for Precision Agriculture." *Computers and Electronics in Agriculture*, 220, 108945.

[3] Mazumder, P., et al. (2024). "Fine-Tuning Pretrained Models for Multi-Crop Plant Disease Identification." *IEEE Access*, 12, 34521-34535.

[4] Alkanan, M., \& Gulzar, Y. (2024). "Efficient Deep Learning Model for Corn Seed Disease Classification Using Transfer Learning." *Agriculture*, 14(7), 1234.

[5] Meghraoui, K., et al. (2024). "A Comprehensive Survey of Machine Learning Approaches for Crop Yield Prediction." *Precision Agriculture*, 25(3), 1890-1920.

[6] Wang, Z., et al. (2024). "Attention-Based LSTM for County-Level Corn Yield Prediction." *Remote Sensing of Environment*, 298, 113812.

[7] Jabed, M. A., \& Murad, S. H. (2024). "CNN-LSTM Hybrid Architecture for Regional Crop Yield Forecasting." *Smart Agricultural Technology*, 8, 100456.

[8] Hiremani, V., et al. (2025). "Federated Deep Learning for Decentralized Crop Disease Identification." *Computers and Electronics in Agriculture*, 225, 109210.

[9] Shaikh, T. A., et al. (2025). "Large Language Models for Agriculture: Opportunities and Challenges." *Artificial Intelligence in Agriculture*, 9, 42-58.

[10] Trentin, C., et al. (2024). "Remote Sensing and Machine Learning for Tree Crop Yield Estimation: A Review." *Frontiers in Plant Science*, 15, 1401875.

[11] Lu, Y., Tan, L., \& Jiang, H. (2021). "Review on Convolutional Neural Network (CNN) Applied to Plant Leaf Disease Classification." *Agriculture*, 11(8), 707.

[12] Saleem, M. H., et al. (2024). "Deep Learning Models for Multi-Crop Leaf Disease Detection." *Plant Methods*, 20, 34.

[13] Radocaj, D., et al. (2024). "A Novel Optimization Module for Transfer Learning in Plant Disease Recognition." *Expert Systems with Applications*, 241, 122456.

[14] Khaki, S., \& Wang, L. (2019). "Crop Yield Prediction Using Deep Neural Networks." *Frontiers in Plant Science*, 10, 621.

[15] De Clercq, M., et al. (2024). "LLMs in Food and Agriculture: Opportunities, Risks, and Challenges." *Trends in Food Science \& Technology*, 145, 104378.

[16] Gujjar, J. P., \& Kumar, H. P. (2025). "Conversational AI Chatbot for Sustainable Agricultural Practices." *IJACSA*, 16(1), 89-97.

[17] Ajith, A., et al. (2025). "A Comprehensive Survey on Machine Learning Approaches for Precision Agriculture." *IEEE Access*, 13, 45678-45702.

[18] Hughes, D. P., \& Salathe, M. (2015). "An Open Access Repository of Images on Plant Health to Enable the Development of Mobile Disease Diagnostics." *arXiv preprint*, arXiv:1511.08060.

[19] Pacal, I., et al. (2024). "A Comprehensive Review of Deep Learning Approaches for Plant Disease Detection." *Applied Soft Computing*, 155, 111459.

[20] Sandler, M., et al. (2018). "MobileNetV2: Inverted Residuals and Linear Bottlenecks." *IEEE CVPR*, pp. 4510-4520.

[21] Hochreiter, S., \& Schmidhuber, J. (1997). "Long Short-Term Memory." *Neural Computation*, 9(8), 1735-1780.

[22] Selvaraju, R. R., et al. (2017). "Grad-CAM: Visual Explanations from Deep Networks via Gradient-based Localization." *IEEE ICCV*, pp. 618-626.

[23] Lundberg, S. M., \& Lee, S.-I. (2017). "A Unified Approach to Interpreting Model Predictions." *NeurIPS*, 30.

[24] IPCC (2019). "2019 Refinement to the 2006 IPCC Guidelines for National Greenhouse Gas Inventories." Volume 4: Agriculture, Forestry and Other Land Use.

[25] Hersbach, H., et al. (2020). "The ERA5 Global Reanalysis." *Quarterly Journal of the Royal Meteorological Society*, 146(730), 1999-2049.

[26] Gorelick, N., et al. (2017). "Google Earth Engine: Planetary-Scale Geospatial Analysis for Everyone." *Remote Sensing of Environment*, 202, 18-27.

[27] Agricultural Census Division (2019). "All India Report on Agriculture Census 2015-16." Ministry of Agriculture and Farmers Welfare, Government of India.

[28] National Statistical Office (2021). "Situation Assessment of Agricultural Households and Land and Livestock Holdings of Households in Rural India, 2019." Ministry of Statistics and Programme Implementation.

[29] Ministry of Finance (2025). "Economic Survey 2024-25." Government of India.

[30] Indian Meteorological Department (2025). "Annual Climate Summary 2024." Ministry of Earth Sciences, Government of India.

\textbf{--- End of Report ---}


\end{document}
